%%%%%%%%%%%%%%%%%%%%%%%%%%%%%%%%%%%%%%%%%%%%%%%%%%%%%%%%%%%%%%%%%%%%%%%%%%%%%%
%  Marketing Management (IM345AT) -- Volume 1 of 5
%  UNIT I : Introduction to Digital Marketing and Content Marketing
%%%%%%%%%%%%%%%%%%%%%%%%%%%%%%%%%%%%%%%%%%%%%%%%%%%%%%%%%%%%%%%%%%%%%%%%%%%%%%

\documentclass{MMTextbook}

\unitnumber{I}
\unittitle{Introduction to Digital Marketing and Content Marketing}
\unitsubtitle{Principles, channels, personas, research tools, and the complete
content marketing engine}
\unithours{07 Hours}
\volumeno{1}

\makeindex[intoc=true,columns=2,title=Index]

\hypersetup{
  pdftitle={Marketing Management IM345AT -- Unit I: Introduction to Digital
             Marketing and Content Marketing},
  pdfsubject={Semester IV, Department of Industrial Engineering and Management,
              RV College of Engineering},
  pdfkeywords={digital marketing, content marketing, buyer persona, SEO,
               IM345AT, marketing management}
}

\begin{document}

%%%%%%%%%%%%%%%%%%%%%%%%%%%%%%%%%%%%%%%%%%%%%%%%%%%%%%%%%%%%%%%%%%%%%%%%%%%%%%
\frontmatter
\pagestyle{mmfront}

\halftitlepage
\maintitlepage

\colophonpage{%
\textbf{\coursetitle\ ---\ The Study Text}\\[2pt]
Volume 1 of 5 \textbullet\ Unit I: Introduction to Digital Marketing and
Content Marketing\\[10pt]
%
Prepared for the course \coursetitle\ (\coursecode), Semester IV, of the
\coursedept, \courseinstitute.\\[10pt]
%
\textbf{Sources.} The syllabus, course outcomes, assessment rubrics and
reference-book list reproduced in the front matter are transcribed from the
official course document. The previous-year questions in Appendix~B are
transcribed from RVCE Continuous Internal Evaluation and Semester End
Examination papers for this course, and where an official scheme of valuation
was released the model answers follow it. The explanatory text is compiled from
classroom lecture notes and, where those were incomplete, from the prescribed
reference texts.\\[10pt]
%
\textbf{Typography.} Set in Nimbus Roman with URW Gothic display headings.
Composed in \LaTeX{} using the \texttt{MMTextbook} class written for this
series. The visual design is derived from \emph{The Legrand Orange Book}
template by Mathias Legrand, with modifications by Vel, distributed through
LaTeXTemplates.com under a Creative Commons BY-NC-SA 4.0 licence; the class
re-implements that design language on standard CTAN packages. All figures are
drawn in \TikZ{} and \texttt{pgfplots}.\\[10pt]
%
\textbf{Status.} Unofficial student-compiled study material. Not issued,
endorsed or verified by the Department or by RV College of Engineering.\\[10pt]
%
2026 edition.}

% ---------------------------------------------------------------- preface
\frontchapter{Preface to Volume 1}

Unit~I is the foundation of \coursetitle, and it carries a structural privilege
that no other unit has: in the Semester End Examination, Question~2 is
compulsory and is always set from Unit~I. A candidate can choose to avoid one
question in each of the other four units. Unit~I cannot be avoided. It is also
the unit from which the first Continuous Internal Evaluation test is drawn
almost in its entirety.

The unit has two halves, and the syllabus names them separately. The first
half, \emph{Introduction to Digital Marketing}, establishes what marketing is,
what makes digital marketing structurally different from what preceded it, and
the practical apparatus of the discipline --- its principles, its channels, the
buyer persona, and the tools used to research competitors and to analyse a
website. The second half, \emph{Content Marketing}, is an engine: an
eight-step cycle that runs from goal setting through creation and distribution
to measurement and improvement, and then begins again.

Chapters~1 to~6 cover the first half; Chapters~7 to~11 cover the second. The
chapter sequence follows the syllabus phrase order exactly, and
Appendix~A records the mapping.

Two editorial decisions in this volume deserve stating plainly.

First, the syllabus descriptor for Unit~I begins at \emph{Introduction to
Digital Marketing} and does not separately name the general
marketing-management foundations --- the mantra of marketing, the entities that
can be marketed, needs and wants and demands, the evolution of company
orientations, the marketing mix. Past papers for this course nevertheless
examine that material inside Unit~I, repeatedly and at both two and ten marks.
It is therefore treated here as the introductory framing of Chapter~1 rather
than as separate content, on the reasoning that one cannot introduce digital
marketing without first fixing what marketing is.

Second, a small number of topics appear in past papers for this course but sit
outside the literal wording of any unit descriptor: the product life cycle,
pricing strategies, and the distinction between market segmentation and the
marketing mix. Rather than silently insert them into a numbered chapter --- which
would breach the rule that the numbered chapters carry the syllabus and nothing
else --- or silently drop them, which would leave real examined material
uncovered, they are set out in Appendix~C and flagged there for what they are.

\vspace{12pt}
\noindent\emph{The whole of the course syllabus is reproduced overleaf. The
unit covered by this volume is Unit~I.}

% ------------------------------------------------- conventions & syllabus
%%%%%%%%%%%%%%%%%%%%%%%%%%%%%%%%%%%%%%%%%%%%%%%%%%%%%%%%%%%%%%%%%%%%%%%%%%%%%%
%  mm-howtouse.tex
%  Editorial conventions for the series.  Shared, identical, across all five
%  volumes so that a reader who learns the apparatus once can use any volume.
%%%%%%%%%%%%%%%%%%%%%%%%%%%%%%%%%%%%%%%%%%%%%%%%%%%%%%%%%%%%%%%%%%%%%%%%%%%%%%

\frontchapter{How to Use This Study Text}

\section*{The shape of a volume}
\addcontentsline{toc}{section}{The shape of a volume}

Each of the five volumes in this series covers exactly one unit of the
IM345AT syllabus and is self-contained. Every volume opens with the official
syllabus for the whole course, so that the unit in hand can always be located
within it, and closes with a set of appendices carrying the concordance,
the question bank, the glossary and a quick-reference sheet.

The numbered chapters of each volume map one-to-one onto the topic phrases of
that unit's official syllabus descriptor. The order of the chapters is the
order of the syllabus. Appendix~A sets out that mapping explicitly, phrase by
phrase, so coverage can be audited rather than assumed.

\section*{The five panel types}
\addcontentsline{toc}{section}{The five panel types}

Material that is not continuous prose is set in one of five panels. Each has a
fixed meaning throughout the series.

\begin{definitionbox}[the term being defined]
A formal definition, worth reproducing close to verbatim in an answer script.
The term appears in the panel heading and again in the index.
\end{definitionbox}

\begin{keybox}[High Yield]
A result, list or distinction that carries disproportionate examination weight.
Where a topic has recurred across papers, the recurrence is recorded as a
bracketed tag. \pyq{CIE-I, 3 Apr 2025; CIE-I, 16 Apr 2026}
\end{keybox}

\begin{exambox}[Model Answer]
An answer written at the length and in the register the marking scheme
rewards. Where an official scheme of valuation was released, the panel follows
it.
\end{exambox}

\begin{examplebox}[Worked Example]
A concrete illustration --- a named company, a named persona, a worked
calculation --- rather than a general statement.
\end{examplebox}

\begin{warnbox}
A common and costly error: a distinction usually collapsed, a term usually
misused, or an answer that looks complete and is not.
\end{warnbox}

A sixth panel, in green, sets out a framework or process as a single
sequence:

\begin{frameworkbox}[Framework]
Step one \ra\ step two \ra\ step three \ra\ back to step one.
\end{frameworkbox}

\section*{How previous-year questions are recorded}
\addcontentsline{toc}{section}{How previous-year questions are recorded}

\begin{keybox}[The Bracket Convention]
A past-paper appearance is recorded \textbf{only} as a bracketed source tag
attached to the material, never as a sentence of narrative. So the text reads
\begin{quote}\sffamily\footnotesize
Push versus pull marketing. \pyq{CIE-I, 3 Apr 2025; CIE-I, 16 Apr 2026; SEE, Jun/Jul 2025 --- 2 M}
\end{quote}
and never ``this was asked in 2025 and again in 2026''. The tag carries the
paper, the date and, where the mark allocation is known, the marks. Multiple
appearances are separated by semicolons.
\end{keybox}

\noindent The abbreviations used inside tags are:

\vspace{4pt}
{\sffamily\small
\begin{tabular}{@{}P{0.14\textwidth}P{0.80\textwidth}@{}}
\toprule[1.1pt]
\rlab{CIE-I} & First Continuous Internal Evaluation test\\
\rlab{CIE-II} & Second Continuous Internal Evaluation test\\
\rlab{CIE-III} & Third Continuous Internal Evaluation test\\
\rlab{SEE} & Semester End Examination\\
\rlab{2 M} & Mark allocation of the question as printed on the paper\\
\bottomrule[1.1pt]
\end{tabular}}

\vspace{10pt}

Appendix~B of each volume reproduces every past-paper question traced to that
unit, in full, grouped by mark weight, each carrying the same bracketed tag.

\section*{Cross-references, figures and the index}
\addcontentsline{toc}{section}{Cross-references, figures and the index}

Cross-references are given as \S\,chapter.section --- so \S\,1.4 is the fourth
section of Chapter~1. Every figure in this series is drawn in \TeX{} rather
than imported as an image, so the diagrams print at full resolution at any
size and use the same typeface as the surrounding text. Figures and tables are
listed after the table of contents.

Terms set in \textbf{\textit{bold italic}} on first use are indexed. The index
at the back of each volume is a real back-of-book index built from those
terms, not a repeat of the contents list.

\section*{Status of this text}
\addcontentsline{toc}{section}{Status of this text}

\begin{warnbox}[Unofficial Document]
This is a student-compiled study text. It is not issued, endorsed or verified
by the Department of Industrial Engineering and Management or by RV College of
Engineering. The syllabus reproduced in the front matter, the assessment
rubrics and the previous-year questions are transcribed from official
documents; the explanatory prose around them is not official. Where the source
lecture notes were incomplete, gaps have been filled from the prescribed
reference texts listed in the front matter. Always cross-check against the
official syllabus and your own class notes.
\end{warnbox}

%%%%%%%%%%%%%%%%%%%%%%%%%%%%%%%%%%%%%%%%%%%%%%%%%%%%%%%%%%%%%%%%%%%%%%%%%%%%%%
%  mm-syllabus.tex
%  Verbatim reproduction of the official IM345AT course document.
%  Shared, byte-identical, across all five volumes of the series.
%%%%%%%%%%%%%%%%%%%%%%%%%%%%%%%%%%%%%%%%%%%%%%%%%%%%%%%%%%%%%%%%%%%%%%%%%%%%%%

\frontchapter{Official Syllabus: IM345AT}

\noindent
{\sffamily\small\color{slatelight}%
Reproduced from the official course document issued by the Department of
Industrial Engineering and Management, RV College of Engineering. The text of
the five unit descriptors below is the authority against which every chapter of
this series is written; nothing outside it appears in the numbered chapters of
any volume.\par}

\vspace{10pt}

\begin{center}
\sffamily\small
\begin{tabular}{@{}P{0.238\textwidth}P{0.238\textwidth}P{0.218\textwidth}P{0.218\textwidth}@{}}
\toprule[1.1pt]
\multicolumn{4}{@{}l}{\bfseries\normalsize Semester: IV \quad\textbullet\quad MARKETING MANAGEMENT}\\
\multicolumn{4}{@{}l}{\color{slatelight}Category: Professional Core Course (Theory)}\\
\midrule
\rlab{Course Code} & IM345AT & \rlab{CIE Marks} & 100 Marks\\
\rlab{Credits (L:T:P)} & 3:0:0 & \rlab{SEE Marks} & 100 Marks\\
\rlab{Total Hours} & 45 L & \rlab{SEE Duration} & 3.00 Hours\\
\bottomrule[1.1pt]
\end{tabular}
\end{center}

\vspace{6pt}

% ---------------------------------------------------------------- unit blocks
\newcommand{\syllunit}[3]{%
  \begin{tcolorbox}[mmbase, colback=white, colframe=ocre!45, boxrule=0.7pt,
      sharp corners, fontupper=\small,
      fonttitle=\sffamily\bfseries\small\color{white},
      coltitle=white, colbacktitle=#3,
      title={#1},
      attach boxed title to top left={xshift=-0.7pt, yshift=-0.7pt},
      boxed title style={sharp corners, boxrule=0pt, colback=#3},
      top=16pt]
    #2
  \end{tcolorbox}}

\syllunit{UNIT--I \hfill 07 Hrs}{%
\textbf{Introduction to Digital Marketing:} Principles of Digital Marketing;
Digital Marketing Channels; Tools to Create Buyer Persona; Competitor Research
Tools, Website Analysis Tools, etc.\par
\vspace{3pt}
\textbf{Content Marketing:} Content Marketing Concepts \& Strategies; Planning,
Creating, Distributing \& Promoting Content; Optimize Website UX \& Landing
Pages; Measure Impact; Metrics \& Performance; Using Content Research for
Opportunities, etc.}{ocre}

\syllunit{UNIT--II \hfill 08 Hrs}{%
\textbf{Social Media Marketing:} Introduction; Major Social Media Platforms for
Marketing; Developing Data-driven Audience \& Campaign Insights; Social Media
for Business; Creation \& Optimization of Social Media Campaigns, etc.\par
\vspace{3pt}
\textbf{Search Engine Optimization:} Search Engine Optimization Fundamentals;
Keywords and SEO Content Plan; SEO \& Business Objectives; Writing SEO Content;
On-site \& off-site SEO; Optimize Organic Search Ranking, etc.}{ocre}

\syllunit{UNIT--III \hfill 07 Hrs}{%
\textbf{Web Analytics \& Google Analytics:} Google Analytics Tools; Web
Analytics Tools, etc.\par
\vspace{3pt}
\textbf{E-mail Marketing:} Effective E-mail Campaigns; E-mail Plan; E-mail
Marketing Campaign Analysis; Measuring Conversions \& keeping up, etc.}{ocre}

\syllunit{UNIT--IV \hfill 07 Hrs}{%
\textbf{Web Design:} Web design, optimization of websites; Publishing a basic
website; User-centered Design and Website Optimization; Design Principles and
Website Copy; Website Metrics \& Developing Insight, etc.\par
\vspace{3pt}
\textbf{Mobile Marketing:} Difference between mobile advertising and marketing,
utilizing mobile marketing for sales promotions, online applications,
etc.}{ocre}

\syllunit{UNIT--V \hfill 07 Hrs}{%
\textbf{Conversion Optimization:} What is AIDAS and its role; website
optimization; what visitors want to see on the website; how to optimize key
element and increase the effect of landing on a particular page\par
\vspace{3pt}
\textbf{Digital Analytics:} Evolution of Digital Analytics, information about
end-to-end customer experience, analyst's influence on business, role as a
change agent, etc.}{ocre}

% ------------------------------------------------------------ course outcomes
\section*{Course Outcomes}
\addcontentsline{toc}{section}{Course Outcomes}

\noindent{\small After completing the course, the students will be able to:\par}
\vspace{4pt}

{\sffamily\small
\begin{longtable}{@{}P{0.07\textwidth}P{0.89\textwidth}@{}}
\toprule[1.1pt]
\rlab{CO1} & Differentiate the benefits drawn by updated marketing mix from
traditional marketing mix for effective marketing management there by to stay
competitive in today's global market-place.\\
\midrule
\rlab{CO2} & Develop an effective holistic marketing atmosphere to efficiently
face the challenges in dynamically changing market.\\
\midrule
\rlab{CO3} & Formulate a potential marketing plan to effectively reach the
targeted market segments, by delivering the value to targeted customers through
practicing sound marketing research.\\
\midrule
\rlab{CO4} & Create new channels to improvise marketing to achieve and maintain
competitive position in globalized market-place.\\
\bottomrule[1.1pt]
\end{longtable}}

% --------------------------------------------------------------- reference books
\section*{Prescribed Reference Books}
\addcontentsline{toc}{section}{Prescribed Reference Books}

{\small
\begin{enumerate}[leftmargin=2.1em]
\item \textit{Marketing Management}, Philip Kotler, Kevin Lane Keller, 15th
      Edition, 2016, Pearson, ISBN: 978-93-325-5718-5.
\item \textit{Digital Marketing --- Strategy, Implementation \& Practice}, Dave
      Chaffey, Fiona Ellis-Chadwick, 7th Edition, 2019, Pearson,
      ISBN: 9781292241623, 1292241624.
\item \textit{Marketing Research}, Donald S. Tull, Del I. Hawkins, 6th Edition,
      Prentice Hall India, ISBN: 8120309618.
\item \textit{Marketing Management --- A South Asian Perspective}, Philip
      Kotler, Kevin Lane Keller, Abraham Koshy, Mithileshwar Jha, 14th Edition,
      2013, Pearson, ISBN: 978-81-317-6716-0.
\item \textit{Marketing Research}, David A. Aaker, V. Kumar, George S. Day, 9th
      Edition, 2008, John Wiley \& Sons, ISBN: 978-265-1791-6.
\end{enumerate}}

% ------------------------------------------------------------------- rubrics
\frontchapter{Assessment Scheme}

\section*{Rubric for the Continuous Internal Evaluation (Theory)}
\addcontentsline{toc}{section}{Rubric for the Continuous Internal Evaluation}

{\sffamily\small
\begin{longtable}{@{}P{0.045\textwidth}P{0.80\textwidth}>{\raggedleft\arraybackslash}P{0.07\textwidth}@{}}
\toprule[1.1pt]
\rowcolor{ocre}\thd{\#} & \thd{Component} & \thd{Marks}\\
\midrule
\rlab{1.} & \textbf{Quizzes:} Quizzes will be conducted in online/offline mode.
TWO QUIZZES will be conducted \& each quiz will be evaluated for 10 marks. The
sum of two quizzes will be the final quiz marks. & \bfseries 20\\
\midrule
\rowcolor{tablealt}
\rlab{2.} & \textbf{Tests:} Students will be evaluated in test, descriptive
questions with different complexity levels (Revised Bloom's Taxonomy Levels:
Remembering, Understanding, Applying, Analyzing, Evaluating and Creating). TWO
tests will be conducted. Each test will be evaluated for 50 marks, adding up to
100 marks. Final test marks will be reduced to 40 marks. & \bfseries 40\\
\midrule
\rlab{3.} & \textbf{Experiential Learning:} Students will be evaluated for their
creativity and practical implementation of the problem. Case study-based
teaching learning (10), Program specific requirements (10), Video based
seminar/presentation/demonstration (20), adding up to 40. & \bfseries 40\\
\midrule
\multicolumn{2}{@{}l}{\bfseries MAXIMUM MARKS FOR THE CIE THEORY} & \bfseries 100\\
\bottomrule[1.1pt]
\end{longtable}}

\section*{Rubric for the Semester End Examination (Theory)}
\addcontentsline{toc}{section}{Rubric for the Semester End Examination}

{\sffamily\small
\begin{longtable}{@{}P{0.12\textwidth}P{0.72\textwidth}>{\raggedleft\arraybackslash}P{0.08\textwidth}@{}}
\toprule[1.1pt]
\rowcolor{ocre}\thd{Q. No.} & \thd{Contents} & \thd{Marks}\\
\midrule
\multicolumn{3}{@{}l}{\bfseries\color{slate} PART A}\\
\midrule
\rlab{1} & Objective type questions covering entire syllabus & \bfseries 20\\
\midrule
\multicolumn{3}{@{}l}{\bfseries\color{slate} PART B \quad\normalfont\itshape(Maximum of TWO sub-divisions only)}\\
\midrule
\rlab{2}      & Unit 1: (Compulsory)      & \bfseries 16\\
\rowcolor{tablealt}
\rlab{3 \& 4} & Unit 2: (Internal Choice) & \bfseries 16\\
\rlab{5 \& 6} & Unit 3: (Internal Choice) & \bfseries 16\\
\rowcolor{tablealt}
\rlab{7 \& 8} & Unit 4: (Internal Choice) & \bfseries 16\\
\rlab{9 \& 10}& Unit 5: (Internal Choice) & \bfseries 16\\
\midrule
\multicolumn{2}{@{}l}{\bfseries TOTAL} & \bfseries 100\\
\bottomrule[1.1pt]
\end{longtable}}

\vspace{8pt}

\begin{keybox}[What the SEE Structure Means for This Volume]
Part A draws twenty marks of objective questions from the \emph{entire}
syllabus, so no unit can be skipped. In Part B, Question~2 is compulsory and is
always set from Unit~I; Units~II to~V each carry a 16-mark internal-choice
pair. A candidate therefore answers Unit~I plus one question from each of the
remaining four units --- five full questions in all.
\end{keybox}


% ---------------------------------------------------------------- contents
\cleardoublepage
\pdfbookmark[0]{Contents}{toc}
\tableofcontents
\cleardoublepage
\listoffigures
\cleardoublepage
\listoftables

%%%%%%%%%%%%%%%%%%%%%%%%%%%%%%%%%%%%%%%%%%%%%%%%%%%%%%%%%%%%%%%%%%%%%%%%%%%%%%
\mainmatter
\pagestyle{mmmain}

%%%%%%%%%%%%%%%%%%%%%%%%%%%%%%%%%%%%%%%%%%%%%%%%%%%%%%%%%%%%%%%%%%%%%%%%%%%%%%
\chapter{Introduction to Digital Marketing}
%%%%%%%%%%%%%%%%%%%%%%%%%%%%%%%%%%%%%%%%%%%%%%%%%%%%%%%%%%%%%%%%%%%%%%%%%%%%%%

\syllabusstrip{Introduction to Digital Marketing.}

\begin{learningoutcomes}
\begin{itemize}
\item What marketing management is, and the mantra that compresses it.
\item The ten entities that can be marketed; needs, wants and demands.
\item How value is created, delivered and judged --- and what satisfaction is.
\item Marketplace, market space and the metamarket.
\item The five company orientations and the four dimensions of holistic marketing\ix{holistic marketing}.
\item What digital marketing is, and why it is not a separate discipline.
\item Traditional marketing, its strengths and its limits, and a costed
      transition to digital.
\item Push marketing against pull marketing.
\item The marketing mix, and how each of the four Ps has evolved in the
      digital era.
\item The technologies that made digital marketing possible.
\item The digital consumer --- the five traits, perceived anonymity, and
      influencers.
\end{itemize}
\end{learningoutcomes}

\section{The Marketing Foundation}

Digital marketing is not a separate discipline from marketing. It is marketing
conducted through digital channels. Every construct in this volume ---
personas, content engines, channel mixes, conversion funnels --- is a digital
expression of something marketing already did. So the foundation has to be
secure before anything else, and it is examined as such.

\subsection{Marketing management defined}

\begin{definitionbox}[Marketing Management]
\keytermas{Marketing management}{marketing management} is the art and science of choosing target
markets and getting, keeping and growing customers through creating, delivering
and communicating superior customer value.
\end{definitionbox}

Three distinct ideas are packed into that sentence, and each carries a
consequence.

\begin{itemize}
\item \sfb{Choosing target markets} is a selection decision. You cannot serve
      everybody, and the attempt to do so produces an offer that fits nobody.
\item \sfb{Getting, keeping and growing customers} is a lifecycle view.
      Acquisition alone is not marketing; a business that acquires well and
      retains badly is buying the same customer twice.
\item \sfb{Creating, delivering and communicating superior value} is the
      actual work --- and note the order. Value is created before it is
      communicated. Communication without created value is advertising a
      promise you cannot keep.
\end{itemize}

\subsection{The mantra of marketing}

\begin{exambox}[Model Answer --- 2 Marks]
\textbf{Create --- Communicate --- Deliver --- Value --- Target --- Profit.}
\pyq{CIE-I, 3 Apr 2025 --- 2 M}

\smallskip
The firm \emph{creates} value, \emph{communicates} it to a chosen
\emph{target}, \emph{delivers} it, and in return earns \emph{profit}.
\end{exambox}

\subsection{Why marketing matters --- three pillars}

\begin{enumerate}
\item \sfb{Financial success.} A firm's financial performance often depends
      directly on its marketing ability. The best operations and the best
      finance function in the industry are worthless if there is insufficient
      demand for what the firm makes.
\item \sfb{Economic engine.} Successful marketing builds demand for products
      and services, and that demand creates jobs.
\item \sfb{Intangible assets.} Marketing builds strong brands and a loyal
      customer base, and these contribute heavily to the firm's overall
      enterprise value --- often more than its physical assets do.
\end{enumerate}

\subsection{Selling is only the tip of the iceberg}

Selling and marketing are not synonyms, and the confusion is expensive.
Selling is the visible tip. Marketing is the mass below the waterline:
research, segmentation, targeting, positioning, product design, pricing,
channel design and communication. There will always be a need for some
selling, but the aim of marketing is to make selling superfluous.

\begin{figure}[htbp]
\centering
\begin{tikzpicture}[font=\sffamily\footnotesize]
  % waterline
  \fill[deepbluepale] (-4.6,0) rectangle (4.6,-4.3);
  \draw[deepblue,thick] (-4.6,0) -- (4.6,0);
  \node[anchor=west,text=deepblue,font=\sffamily\scriptsize\itshape]
    at (-4.55,0.22) {waterline: what the customer sees};
  % tip
  \fill[ocre] (0,1.55) -- (-0.95,0) -- (0.95,0) -- cycle;
  \node[anchor=south,text=slate,font=\sffamily\bfseries\small] at (0,1.72) {SELLING};
  % submerged mass
  \fill[ocre!35] (-0.95,0) -- (0.95,0) -- (3.7,-4.1) -- (-3.7,-4.1) -- cycle;
  \node[text=slate,font=\sffamily\bfseries\small] at (0,-0.62) {MARKETING};
  \node[text=slate,align=center,font=\sffamily\scriptsize] at (0,-1.75)
    {consumer research \ \textbullet\ \ product development\\
     segmentation \ \textbullet\ \ targeting \ \textbullet\ \ positioning\\
     pricing \ \textbullet\ \ channel design\\
     communication \ \textbullet\ \ service};
\end{tikzpicture}
\caption[Selling as the tip of the marketing iceberg]%
        {Selling is the visible tip; marketing is the submerged mass that makes
         the tip possible. The aim of the work below the waterline is to reduce
         how much work the tip has to do.}
\label{fig:iceberg}
\end{figure}

\begin{quotebox}[Peter Drucker]
The aim of marketing is to make selling superfluous. The aim is to know and
understand the customer so well that the product or service fits him and sells
itself. Ideally marketing should result in a customer who is ready to buy. All
that should be needed then is to make the product or service available.
\end{quotebox}

\subsection{The ten entities that can be marketed}

Marketing is not confined to physical products. Ten classes of thing can be
marketed, and the official scheme for this course lists exactly ten.
\pyq{CIE-I, 3 Apr 2025 --- 10 M}

\begin{mmlongtable}{The ten entities that can be marketed}{tab:entities}%
{@{}P{0.15\textwidth}P{0.42\textwidth}P{0.36\textwidth}@{}}
\rowcolor{ocre}\thd{Entity} & \thd{What is marketed} & \thd{Example}\\
\midrule
\endfirsthead
\rowcolor{ocre}\thd{Entity} & \thd{What is marketed} & \thd{Example}\\
\midrule
\endhead
\rlab{Goods} & Physical, tangible products --- the bulk of most nations'
production & Cars, smartphones, packaged food\\
\rowcolor{tablealt}
\rlab{Services} & Intangible performances that cannot be stored &
Airlines, banking, hotels, salons, consultants\\
\rlab{Events} & Time-based occasions promoted to an audience &
Olympics, trade fairs, IPL, company anniversaries\\
\rowcolor{tablealt}
\rlab{Experiences} & Orchestrated combinations of goods and services &
Disneyland, adventure treks, escape rooms\\
\rlab{Persons} & Celebrity and personal branding &
Sportspersons, actors, CEOs, political leaders\\
\rowcolor{tablealt}
\rlab{Places} & Cities, states, regions and nations competing for tourists and
investment & ``Incredible India'', Singapore, Dubai\\
\rlab{Properties} & Intangible rights of ownership, real or financial &
Real estate, stocks, bonds\\
\rowcolor{tablealt}
\rlab{Organizations} & The institution itself as a brand &
Universities, NGOs, corporate brand campaigns\\
\rlab{Information} & Production, packaging and distribution of knowledge &
Encyclopaedias, magazines, schools, market research\\
\rowcolor{tablealt}
\rlab{Ideas} & Social causes and concepts &
``Swachh Bharat'', anti-smoking and road-safety campaigns\\
\end{mmlongtable}

\subsection{Needs, wants and demands}

\begin{definitionbox}[Needs, Wants, Demands]
\keytermas{Needs}{needs} are the basic human requirements --- air, food, water,
clothing, shelter, and at a higher level recreation, education and
entertainment. Needs are not created by marketers.

\keytermas{Wants}{wants} are needs directed at specific objects that might satisfy
them. Wants are shaped by society, culture and individual personality.

\keytermas{Demands}{demands} are wants for a specific product backed by an ability and
willingness to pay.
\end{definitionbox}

An Indian consumer \emph{needs} food, \emph{wants} a masala dosa, and
\emph{demands} one only if there is money in the pocket. Many people want a
Mercedes; only a few can demand one. The distinction matters because it locates
the marketer's actual leverage: needs cannot be manufactured, wants can be
shaped, and demand can be enabled through pricing, credit and distribution.

\begin{exambox}[Core Concepts of Marketing --- 2 Marks]
\pyq{CIE-I, 16 Apr 2026 --- 2 M} Acceptable answers are any two of: needs,
wants, demand, value, exchange, markets. Extended lists also admit target
markets, segmentation, positioning, offerings, brands, satisfaction, supply
chain, competition and the marketing environment.
\end{exambox}

\section{The Value Creation Engine}

Value moves through the organisation in three connected steps: decide the
audience, build the solution, and measure the result.

\begin{figure}[htbp]
\centering
\begin{tikzpicture}[font=\sffamily\footnotesize,
    stage/.style={draw=ocre,line width=0.9pt,fill=ocrepale,rounded corners=4pt,
                  minimum width=3.5cm,minimum height=1.05cm,align=center,
                  font=\sffamily\bfseries\small,anchor=center},
    item/.style={align=center,font=\sffamily\scriptsize,text=slate,
                 text width=3.7cm,anchor=north},
    arr/.style={-{Latex[length=2.6mm]},ocre,line width=1pt}]

  % --- stage boxes on the top row
  \node[stage] (a) at (0,0)    {1. THE AUDIENCE};
  \node[stage] (b) at (5.3,0)  {2. THE SOLUTION};
  \node[stage] (c) at (10.6,0) {3. THE RESULT};

  \draw[arr] (a) -- (b);
  \draw[arr] (b) -- (c);

  % --- descriptive blocks, hung from a fixed y so heights cannot interact
  \node[item] at (0,-0.95)    {%
    target markets\\ segmentation\\ positioning\\[3pt]
    \textit{who to serve, and where to stand in their mind}};
  \node[item] at (5.3,-0.95)  {%
    value proposition\\ the offering\\ the brand\\[3pt]
    \textit{the benefits, the bundle that delivers them, and the known source}};
  \node[item] at (10.6,-0.95) {%
    value: quality, service, price\\ satisfaction\\[3pt]
    \textit{the QSP triad, and the customer's judgement of it}};

  % --- feedback path, routed outside the columns so it crosses no text
  \draw[arr,dashed]
     (10.6,-3.75) -- (10.6,-4.5) -- (-2.35,-4.5) -- (-2.35,0) -- (a.west);
  \node[anchor=south,font=\sffamily\scriptsize\itshape,text=ocredark]
     at (4.4,-4.46) {satisfaction data re-enters the definition of the audience};
\end{tikzpicture}
\caption[The value creation engine]%
        {The value creation engine. Value is defined by the audience chosen,
         built into the solution, and judged in the result --- and the judgement
         feeds back into the next definition of the audience.}
\label{fig:valueengine}
\end{figure}

\begin{enumerate}
\item \sfb{The audience} --- target markets, segmentation and positioning.
      Identifying who to serve and where to stand in their mind.
\item \sfb{The solution} --- the \keyterm{value proposition} (the set of
      benefits offered to satisfy the need), the \keyterm{offering} (the
      combination of products, services, information and experiences that
      delivers it), and the \keyterm{brand} (an offering from a known source).
\item \sfb{The result} --- value (the combination of quality, service and
      price, the \emph{QSP triad}) and \keyterm{satisfaction} (a person's
      judgement of a product's perceived performance compared with
      expectations).
\end{enumerate}

\subsection{Satisfaction, and the difference between satisfied and delighted}

\begin{keybox}[The Satisfaction Rule]
If performance $<$ expectations, the customer is \textbf{dissatisfied}.\\
If performance $=$ expectations, the customer is \textbf{satisfied}.\\
If performance $>$ expectations, the customer is \textbf{delighted}.

\smallskip
Delight, not mere satisfaction, is what produces advocacy\ix{advocacy}. A satisfied customer
got what was promised and owes you nothing; a delighted customer got more than
was promised and tells people so. This asymmetry is why the whole of earned
media (\S\,3.2) rests on exceeding expectations rather than meeting them.
\end{keybox}

\begin{figure}[htbp]
\centering
\begin{tikzpicture}[font=\sffamily\footnotesize]
  % direction-of-travel arrow
  \draw[-{Latex[length=2.2mm]},ocre,line width=1pt] (0.4,1.75) -- (10.8,1.75);
  \node[anchor=south,text=ocre,font=\sffamily\scriptsize] at (5.6,1.8)
    {increasing perceived performance against a fixed expectation};

  \draw[rulegrey,line width=0.8pt] (0,0) -- (11.2,0);

  % Labels are stacked on two lines so that adjacent bands cannot collide.
  \foreach \x/\lab/\col/\rel in {%
     1.7/DISSATISFIED/warnred/{$<$},
     5.6/SATISFIED/slatelight/{$=$},
     9.5/DELIGHTED/greenok/{$>$}}{
    \fill[\col] (\x,0) circle (0.15);
    \node[anchor=south,text=\col,font=\sffamily\bfseries\small] at (\x,0.32) {\lab};
    \node[anchor=north,text=slate,font=\sffamily\scriptsize,align=center,
          text width=2.9cm] at (\x,-0.26)
      {performance\\ \rel\ expectations};
  }
  \node[anchor=north,text=warnred,font=\sffamily\scriptsize\itshape,
        align=center,text width=2.9cm]   at (1.7,-1.3) {churn,\\ complaint};
  \node[anchor=north,text=slatelight,font=\sffamily\scriptsize\itshape,
        align=center,text width=2.9cm]   at (5.6,-1.3) {repeat\\ purchase};
  \node[anchor=north,text=greenok,font=\sffamily\scriptsize\itshape,
        align=center,text width=2.9cm]   at (9.5,-1.3) {advocacy,\\ referral};
\end{tikzpicture}
\caption[The satisfaction scale]%
        {The satisfaction scale. Only the third band generates earned media;
         the second merely avoids losing the customer.}
\label{fig:satisfaction}
\end{figure}

\subsection{Marketplace, market space and the metamarket}

\begin{definitionbox}[Marketplace and Market Space]
A \keyterm{marketplace} is a physical or digital location where buyers and
sellers come together to exchange goods and services. It is concrete --- a
mall, a mandi, a showroom.

A \keyterm{market space} is a conceptual or virtual environment in which
commercial value is created through digital networks and communication
technologies. It is abstract and has no geography --- Amazon, an app, a
marketplace platform.
\end{definitionbox}

\begin{keybox}[High Yield]
One-line contrast for a two-mark answer: \textbf{marketplace = physical;
market space = digital.}
\pyq{CIE-I, 3 Apr 2025; CIE-I, 16 Apr 2026 --- 2 M each}
\end{keybox}

A third, related term completes the set. A \keyterm{metamarket} is a cluster of
complementary products and services that are closely related in the mind of the
consumer but spread across different industries. The automobile metamarket
contains manufacturers, dealers, insurers, financiers, service centres and auto
magazines --- none of which is in the same industry, all of which are in the
same purchase.

\section{Company Orientations and Holistic Marketing}

\subsection{The evolution of company orientations}

\begin{mmlongtable}{The five company orientations}{tab:orientations}%
{@{}P{0.20\textwidth}P{0.44\textwidth}P{0.29\textwidth}@{}}
\rowcolor{ocre}\thd{Orientation} & \thd{Central belief} & \thd{Focus}\\
\midrule
\endfirsthead
\rowcolor{ocre}\thd{Orientation} & \thd{Central belief} & \thd{Focus}\\
\midrule
\endhead
\rlab{Production concept} & Consumers prefer products that are widely available
and inexpensive & Manufacturing efficiency, mass distribution, low cost\\
\rowcolor{tablealt}
\rlab{Product concept} & Consumers favour products offering the most quality,
performance or innovative features & Continuous product improvement\\
\rlab{Selling concept} & Consumers will not buy enough unless prodded &
Aggressive promotion and a large sales force --- the tip of the iceberg\\
\rowcolor{tablealt}
\rlab{Marketing concept} & Find the right products for your customers, not the
right customers for your products & Customer needs, integrated marketing,
profitability through satisfaction\\
\rlab{Holistic marketing} & Everything matters in marketing, and a broad,
integrated perspective is necessary & Internal, integrated, performance and
relationship marketing\\
\end{mmlongtable}

\begin{warnbox}[Do Not Collapse the Middle Two]
The product concept and the selling concept are routinely confused. The product
concept believes a \emph{better} product will sell itself, and fails through
engineering-led blindness to what customers actually want. The selling concept
believes \emph{any} product will sell if pushed hard enough, and fails through
churn\ix{churn} --- it can move the first unit but never the second.
\end{warnbox}

\subsection{Holistic marketing --- four dimensions}

\begin{itemize}
\item \sfb{Internal marketing} --- aligning the marketing department, senior
      management and every other department so that the whole firm embraces
      marketing principles.
\item \sfb{Integrated marketing} --- designing communications, products and
      services, channels and price so that they reinforce rather than
      contradict one another.
\item \sfb{Performance marketing} --- measuring returns: sales revenue, brand
      equity and customer equity\ix{customer equity}, plus the ethical, environmental, legal and
      social effects of marketing activity.
\item \sfb{Relationship marketing} --- building deep, enduring relationships
      with the four constituents who sustain the business: customers,
      employees, marketing partners and the financial community.
\end{itemize}

\subsection{The eight marketing management tasks}

\begin{frameworkbox}[The Eight Tasks]
1. Developing market strategies and plans \ra{} 2. Capturing marketing insights
\ra{} 3. Connecting with customers \ra{} 4. Building strong brands \ra{}
5. Creating value \ra{} 6. Delivering value \ra{} 7. Communicating value \ra{}
8. Creating successful long-term growth --- which loops back to task 1.
\end{frameworkbox}

\subsection{The market ecosystem and the profit equation}

An \emph{industry} is a collection of sellers; a \emph{market} is a collection
of buyers. Sellers send communication and goods and services to the market;
buyers send back money and information. The whole loop resolves into one
equation.

\begin{figure}[htbp]
\centering
\begin{tikzpicture}[font=\sffamily\footnotesize,
    blk/.style={draw=slate,line width=0.9pt,fill=white,rounded corners=4pt,
                minimum width=3.2cm,minimum height=1.2cm,align=center,
                font=\sffamily\bfseries\small,text=slate}]
  \node[blk] (ind) at (0,0) {INDUSTRY\\[-1pt]{\scriptsize\mdseries a collection of sellers}};
  \node[blk] (mkt) at (7.4,0) {MARKET\\[-1pt]{\scriptsize\mdseries a collection of buyers}};

  \draw[-{Latex[length=2.6mm]},ocre,line width=1pt]
    ($(ind.north east)+(0,-0.18)$) to[bend left=22]
    node[above,midway,text=ocre,font=\sffamily\scriptsize]{goods and services}
    ($(mkt.north west)+(0,-0.18)$);
  \draw[-{Latex[length=2.6mm]},ocre,line width=1pt]
    ($(ind.east)+(0,0.02)$) to[bend left=8]
    node[above,midway,text=ocre,font=\sffamily\scriptsize]{communication}
    ($(mkt.west)+(0,0.02)$);
  \draw[-{Latex[length=2.6mm]},deepblue,line width=1pt]
    ($(mkt.south west)+(0,0.18)$) to[bend left=22]
    node[below,midway,text=deepblue,font=\sffamily\scriptsize]{money}
    ($(ind.south east)+(0,0.18)$);
  \draw[-{Latex[length=2.6mm]},deepblue,line width=1pt]
    ($(mkt.south west)+(0,0.55)$) to[bend left=8]
    node[below,midway,text=deepblue,font=\sffamily\scriptsize,yshift=-1pt]{information}
    ($(ind.south east)+(0,0.55)$);
\end{tikzpicture}
\caption[The industry--market exchange loop]%
        {The exchange loop. Two flows out, two flows back. Marketing manages
         all four.}
\label{fig:exchangeloop}
\end{figure}

\begin{keybox}[The Profit Equation]
\[ \text{Profit} \;=\; \text{Revenues} \;-\; \text{Costs} \]
Every marketing decision either raises revenue --- more customers, a higher
price, more purchases per customer --- or lowers cost --- cheaper acquisition,
better retention, more efficient channels. If it does neither, it is not
marketing. It is decoration.
\end{keybox}

Two further constructs complete the ecosystem.

\begin{description}
\item[Supply chain] The longer channel stretching from raw materials through
components to finished products carried to final buyers. For coffee: Ethiopian
farmers grow and harvest, then sell to a Fair Trade cooperative; the beans are
washed, dried and packaged; an Alternative Trading Organisation transports them;
and they are sold through retail channels. Each link captures a share of the
total value.
\item[Competition] All the actual and potential rival offerings and substitutes
a buyer might consider --- not merely direct competitors, but anything
competing for the same need and the same rupee.
\end{description}

\section{What Digital Marketing Is}

\begin{definitionbox}[Digital Marketing]
\keytermas{Digital marketing}{digital marketing} is the well-targeted, conversion-oriented,
quantifiable and interactive marketing of products or services using digital
technologies, in order to reach customers, convert them into clients and retain
them in a sustainable fashion.
\end{definitionbox}

Four words in that definition do the work. \emph{Well-targeted}: the audience
can be specified rather than approximated. \emph{Conversion-oriented}: the
objective is a defined action, not an impression\ix{impression}. \emph{Quantifiable}: every
step leaves a data trail. \emph{Interactive}: the audience can answer back.

Through digital channels a business does not merely share its products and
services online. It can attract prospects, engage them, and convert them into
paying customers, while measuring every step. The speed and simplicity with
which digital media transmits information and supports a business is what makes
it structurally different from everything that came before.

The world is now super-connected, and marketing and advertising are no longer
what they once were. This is particularly true because of the rise of social
networking, which changed how organisations speak with potential and existing
customers. Digital marketing is an aggregate term, used where advertising and
marketing meet web technology and online media platforms.

\begin{keybox}[The Central Idea of This Course]
``Digital marketing isn't about technology; it's about people.'' Technology is
only interesting, from a marketing perspective, when it connects people with
other people more effectively. Understanding the underlying code is secondary.
Understanding the human being on the other side of the screen is the job.
\end{keybox}

\section{Traditional Marketing and the Transition to Digital}

\subsection{Traditional marketing}

Traditional marketing lets businesses promote their products or services on
print media, radio and television commercials, billboards, business cards and
numerous similar channels where the internet and social platforms are not used.
It is, in effect, everything except digital means of branding a product.

\begin{itemize}
\item Traditional approaches had constrained reach and limited scope to
      influence buying behaviour, and they were not quantifiable.
\item Tangible forms include business cards, print advertisements in newspapers
      and magazines, posters, brochures, television and radio commercials, and
      billboards.
\item An often overlooked form is \sfb{referral and networking} --- people find
      a business through a referral and a rapport is built over time.
\item Because of its longevity, people are accustomed to it. Reading a
      billboard or spotting an advertisement in a magazine remain familiar
      everyday activities.
\item It usually reaches a local audience, though it is not strictly limited to
      one.
\end{itemize}

\begin{examplebox}[Where Traditional Marketing Still Wins]
Targeting a genuinely local audience; materials the customer physically keeps
and returns to; a more personal feel in a high-trust sale; and a simpler
process with fewer moving parts and no platform dependency.
\end{examplebox}

Its primary weaknesses are the mirror image. Results are not easily measured
and often cannot be measured at all. It is generally more costly. And it is
static --- there is no way to interact with the audience. You are throwing
information in front of people and hoping they act.

\subsection{Traditional marketing against digital marketing}

\begin{keybox}[High Yield]
This comparison is examined at both ten and eight marks.
Learn at least seven rows and keep one example per row.
\pyq{CIE-I, 3 Apr 2025 --- 10 M; CIE-I, 16 Apr 2026 --- 10 M; SEE, Oct/Nov 2023, Q2a --- 8 M}
\end{keybox}

\begin{mmlongtable}{Traditional marketing against digital marketing}{tab:tradvsdig}%
{@{}P{0.17\textwidth}P{0.38\textwidth}P{0.38\textwidth}@{}}
\rowcolor{ocre}\thd{Aspect} & \thd{Traditional Marketing} & \thd{Digital Marketing}\\
\midrule
\endfirsthead
\rowcolor{ocre}\thd{Aspect} & \thd{Traditional Marketing} & \thd{Digital Marketing}\\
\midrule
\endhead
\rlab{Communication} & Unidirectional --- the organisation talks at its
audience & Bidirectional --- customers ask questions, complain and make
suggestions back\\
\rowcolor{tablealt}
\rlab{Medium} & Phone calls, letters, print, television, radio, billboards &
Social media, websites, chat, apps, e-mail, SEO, PPC\\
\rlab{Reach} & Local, or mass depending on medium; geographically restricted &
Global and theoretically infinite, yet laser-targeted\\
\rowcolor{tablealt}
\rlab{Cost} & Often expensive --- television, print production, media buying &
Cost-effective --- pay-per-click and organic reach options\\
\rlab{Targeting} & Broad and imprecise & Precise --- age, location, interests,
behaviour, intent\\
\rowcolor{tablealt}
\rlab{Interactivity} & One-way communication & Two-way, real-time engagement\\
\rlab{Measurability} & Hard to track ROI accurately; largely unquantifiable &
Easily trackable with analytics and KPIs, in real time\\
\rowcolor{tablealt}
\rlab{Flexibility} & Less flexible once launched; rigid, long-term campaigns &
Highly flexible; adjustable in real time\\
\rlab{Time to launch} & Requires long planning and production cycles &
Quick to launch, refine and adapt\\
\rowcolor{tablealt}
\rlab{Direction of pull} & Outbound --- push & Inbound --- the goal is for
people to find you\\
\end{mmlongtable}

\subsection{A strategy for transitioning from traditional to digital}

\begin{exambox}[Model Answer --- 10 Marks]
\pyq{CIE-I, 3 Apr 2025; CIE-I, 16 Apr 2026 --- 10 M}
A complete answer states the benefits and limitations of both approaches ---
use Table~\ref{tab:tradvsdig} --- and then commits to a sequenced transition.

\begin{enumerate}
\item \sfb{Start with a hybrid approach.} Run both in parallel initially rather
      than cutting traditional spend overnight. Cutting first destroys known
      demand before the replacement exists.
\item \sfb{Build the owned foundation.} A quality company website and a social
      media presence, because every later step needs somewhere to send traffic.
\item \sfb{Layer on SEO and content marketing.} Blogs, videos and landing
      pages, to build compounding organic traffic.
\item \sfb{Run targeted paid advertising.} Google Ads and social advertising to
      buy immediate reach while the organic asset builds.
\item \sfb{Leverage e-mail marketing and influencers} for retention and
      borrowed credibility respectively.
\item \sfb{Track performance with analytics tools}, defining the KPIs
      \emph{before} launch rather than after.
\item \sfb{Gradually reduce traditional spending} based on measured results,
      not on assumption.
\item \sfb{Upskill or hire.} Train the existing team or bring in a dedicated
      digital marketing expert.
\end{enumerate}

\smallskip
\textbf{Limitations to name for digital}, so the answer is balanced: it
requires technical skills; online competition is high; there are privacy and
data-protection concerns; and it depends on internet access.
\end{exambox}

\section{Push and Pull Marketing}

\begin{keybox}[The Single Most Repeated Two-Mark Question]
Push against pull is the most repeated short question in this course. Learn the definitions close to verbatim and keep two examples ready.
\pyq{CIE-I, 3 Apr 2025; CIE-I, 16 Apr 2026; SEE, Jun/Jul 2025 --- 2 M each}
\end{keybox}

\begin{figure}[htbp]
\centering
\begin{tikzpicture}[font=\sffamily\footnotesize,
    nd/.style={draw=slate!55,fill=white,rounded corners=3pt,
               minimum width=2.05cm,minimum height=0.78cm,align=center,
               font=\sffamily\scriptsize\bfseries,text=slate}]

  % ---- push row
  \node[text=ocre,font=\sffamily\bfseries\small,anchor=west] at (-0.2,1.5) {PUSH};
  \node[nd] (p1) at (1.9,1.5) {PRODUCER};
  \node[nd] (p2) at (4.7,1.5) {CHANNEL};
  \node[nd] (p3) at (7.5,1.5) {CONSUMER};
  \draw[-{Latex[length=2.6mm]},ocre,line width=1.1pt] (p1)--(p2);
  \draw[-{Latex[length=2.6mm]},ocre,line width=1.1pt] (p2)--(p3);
  \node[text=slatelight,font=\sffamily\scriptsize,anchor=west] at (8.75,1.5)
    {consumer is \textit{passive}};

  % ---- pull row
  \node[text=deepblue,font=\sffamily\bfseries\small,anchor=west] at (-0.2,-0.2) {PULL};
  \node[nd] (q1) at (1.9,-0.2) {PRODUCER};
  \node[nd] (q2) at (4.7,-0.2) {CHANNEL};
  \node[nd] (q3) at (7.5,-0.2) {CONSUMER};
  \draw[{Latex[length=2.6mm]}-,deepblue,line width=1.1pt] (q1)--(q2);
  \draw[{Latex[length=2.6mm]}-,deepblue,line width=1.1pt] (q2)--(q3);
  \node[text=slatelight,font=\sffamily\scriptsize,anchor=west] at (8.75,-0.2)
    {consumer is \textit{searching}};

  \node[text=slate,font=\sffamily\scriptsize,anchor=west] at (0.35,0.65)
    {\textit{demand is created by the seller}};
  \node[text=slate,font=\sffamily\scriptsize,anchor=west] at (0.35,-1.05)
    {\textit{demand is created by the buyer}};
\end{tikzpicture}
\caption[Push and pull marketing]%
        {Push and pull differ in the direction of initiative, and therefore in
         who creates the demand.}
\label{fig:pushpull}
\end{figure}

\begin{mmlongtable}{Push marketing against pull marketing}{tab:pushpull}%
{@{}P{0.19\textwidth}P{0.37\textwidth}P{0.37\textwidth}@{}}
\rowcolor{ocre}\thd{Basis} & \thd{Push Marketing} & \thd{Pull Marketing}\\
\midrule
\endfirsthead
\rowcolor{ocre}\thd{Basis} & \thd{Push Marketing} & \thd{Pull Marketing}\\
\midrule
\endhead
\rlab{Definition} & The brand pushes the product towards the customer &
The brand encourages customers to actively seek it out\\
\rowcolor{tablealt}
\rlab{Direction} & Producer \ra{} channel \ra{} consumer &
Consumer \ra{} producer\\
\rlab{Typical tools} & Television and print advertising, trade discounts,
dealer incentives, direct selling, telemarketing, point-of-sale displays &
SEO, content marketing, social media, word of mouth, brand building, public
relations\\
\rowcolor{tablealt}
\rlab{Consumer state} & Passive; unaware of the need &
Active; already searching\\
\rlab{Demand created by} & The seller & The buyer\\
\rowcolor{tablealt}
\rlab{Cost pattern} & High recurring spend; stops when spend stops &
High upfront effort; compounds over time\\
\rlab{Best for} & New products, impulse goods, low-involvement purchases &
Established brands, high-involvement and considered purchases\\
\end{mmlongtable}

\section{The Marketing Mix in the Digital Era}

\subsection{The four Ps}

\begin{definitionbox}[Marketing Mix]
The \keyterm{marketing mix} is the set of controllable tactical tools that a
firm blends to produce the response it wants from its chosen target market\ix{target market}:
Product, Price, Place and Promotion.
\end{definitionbox}

\begin{figure}[htbp]
\centering
\begin{tikzpicture}[font=\sffamily\footnotesize,
    p/.style={draw=ocre,line width=1pt,fill=ocrepale,rounded corners=5pt,
              text width=3.1cm,align=center,inner sep=6pt,text=slate,
              font=\sffamily\scriptsize},
    hd/.style={font=\sffamily\bfseries\small,text=ocredark}]

  \node[draw=slate,line width=1pt,fill=white,circle,minimum size=2.15cm,
        align=center,font=\sffamily\bfseries\small,text=slate] (c) at (0,0)
        {TARGET\\MARKET};

  \node[p] (prod) at (-4.3,2.0)  {\textbf{PRODUCT}\\[2pt]
      variety, quality, design, features, brand name, packaging, sizes,
      services, warranties, returns};
  \node[p] (pric) at (4.3,2.0)   {\textbf{PRICE}\\[2pt]
      list price, discounts, allowances, payment period, credit terms};
  \node[p] (plac) at (-4.3,-2.0) {\textbf{PLACE}\\[2pt]
      channels, coverage, assortments, locations, inventory, transport};
  \node[p] (prom) at (4.3,-2.0)  {\textbf{PROMOTION}\\[2pt]
      sales promotion, advertising, sales force, public relations, direct
      marketing};

  \foreach \n in {prod,pric,plac,prom}{
     \draw[-{Latex[length=2.2mm]},ocre,line width=0.9pt] (\n) -- (c);}
\end{tikzpicture}
\caption[The marketing mix]%
        {The four Ps are the controllable levers; the target market is what
         they are aimed at. Changing the target changes every lever.}
\label{fig:fourps}
\end{figure}

\subsection{How the four Ps have evolved in the digital era}

\begin{mmlongtable}{The digital evolution of the four Ps}{tab:digitalps}%
{@{}P{0.11\textwidth}P{0.46\textwidth}P{0.35\textwidth}@{}}
\rowcolor{ocre}\thd{P} & \thd{Digital evolution} & \thd{Company example}\\
\midrule
\endfirsthead
\rowcolor{ocre}\thd{P} & \thd{Digital evolution} & \thd{Company example}\\
\midrule
\endhead
\rlab{Product} & Digital and personalised products; customisable, data-driven,
and increasingly delivered in digital formats --- software, online services,
virtual goods & Spotify uses listening data to generate personalised playlists,
making the product itself different for every user\\
\rowcolor{tablealt}
\rlab{Price} & Dynamic and transparent pricing; real-time price changes,
personalised offers, and instant comparison by the customer &
Amazon adjusts prices dynamically using demand, competition and user behaviour\\
\rlab{Place} & Online channels and omnichannel presence; a shift from physical
stores to e-commerce, apps and marketplaces with a seamless online--offline
experience & Nike links its website, mobile app and physical stores through
click-and-collect and a store locator\\
\rowcolor{tablealt}
\rlab{Promotion} & Digital and interactive marketing; a shift from mass media
to targeted digital campaigns, social media, influencer and content marketing &
Coca-Cola's ``Share a Coke'' campaign ran as interactive, personalised,
shareable digital content\\
\end{mmlongtable}

\begin{exambox}[Examined As]
``How have the 4Ps of marketing evolved in the digital era? Provide examples of
companies effectively leveraging digital strategies for each element.''
\pyq{CIE-I, 3 Apr 2025 --- 10 M}
Answer with Table~\ref{tab:digitalps}, one named company per P, and close by
naming the service extensions below.
\end{exambox}

\subsection{The seven Ps and the four Cs}

Two standard extensions complete the framework.

\begin{itemize}
\item The \sfb{seven Ps of services marketing} add three service-specific
      levers to the original four: \sfb{People} (everyone who touches the
      customer), \sfb{Process} (the sequence the customer is put through) and
      \sfb{Physical evidence} (the tangible cues that make an intangible
      service credible).
\item The \sfb{four Cs} are the customer-side mirror of the four Ps:
      \sfb{Customer solution} for Product, \sfb{Customer cost} for Price,
      \sfb{Convenience} for Place and \sfb{Communication} for Promotion. The
      relabelling is not cosmetic --- it forces each lever to be stated from
      the buyer's point of view rather than the seller's.
\end{itemize}

\section{The Technology Behind Digital Marketing}

Developments in technology and the evolution of marketing are inextricably
intertwined, and the adoption pattern is always the same.

\begin{figure}[htbp]
\centering
\begin{tikzpicture}[font=\sffamily\footnotesize,
    st/.style={draw=ocre,line width=0.9pt,fill=white,
               chamfered rectangle,chamfered rectangle xsep=0.5cm,
               text width=2.35cm,align=center,inner sep=4pt,
               font=\sffamily\scriptsize,text=slate}]
  \node[st,fill=ocrepale] (s1) at (0,0)    {\textbf{1. Emergence}\\ confined to
      technologists and early adopters};
  \node[st,fill=ocrepale] (s2) at (3.65,0) {\textbf{2. Foothold}\\ becomes
      popular; appears on the marketing radar};
  \node[st,fill=ocrepale] (s3) at (7.3,0)  {\textbf{3. Exploration}\\ innovative
      marketers harness it to reach an audience};
  \node[st,fill=ocrepale] (s4) at (10.95,0){\textbf{4. Mainstream}\\ absorbed
      into standard marketing practice};
  \foreach \a/\b in {s1/s2,s2/s3,s3/s4}{
    \draw[-{Latex[length=2.4mm]},ocre,line width=1pt] (\a)--(\b);}
\end{tikzpicture}
\caption[The technology adoption pattern]%
        {The four-step adoption pattern. The printing press, radio, television
         and the internet each ran the same course.}
\label{fig:adoption}
\end{figure}

The printing press, radio, television and now the internet are all
breakthroughs that permanently altered the relationship between marketers and
consumers, on a global scale.

\subsection{Timeline --- from Sputnik to the World Wide Web}

\begin{mmlongtable}{The enabling technologies of digital marketing}{tab:timeline}%
{@{}P{0.11\textwidth}P{0.46\textwidth}P{0.35\textwidth}@{}}
\rowcolor{ocre}\thd{Year} & \thd{Milestone} & \thd{Why it matters to marketers}\\
\midrule
\endfirsthead
\rowcolor{ocre}\thd{Year} & \thd{Milestone} & \thd{Why it matters to marketers}\\
\midrule
\endhead
\rlab{1957} & The USSR launches Sputnik & Signals that the United States is
falling behind; triggers heavy US investment in science and technology\\
\rowcolor{tablealt}
\rlab{1958} & The US Department of Defense sets up ARPA, the Advanced Research
Projects Agency & The institutional parent of the internet\\
\rlab{early 1960s} & E-mail exists only within single mainframes & The killer
application is born, but cannot yet cross machines\\
\rowcolor{tablealt}
\rlab{1971} & Ray Tomlinson writes the first program to send mail between host
computers on ARPANET, choosing the \texttt{@} symbol to separate user from host
& E-mail becomes a network application --- as a programmer's afterthought\\
\rlab{1974} & Vinton Cerf, the ``father of the internet'', and Robert Kahn
define the TCP protocol at DARPA & Different networks can now be
interconnected\\
\rowcolor{tablealt}
\rlab{1983} & ARPANET adopts TCP/IP; the domain name system is invented
(\texttt{.com}, \texttt{.net}) & Widely considered the true birth of the
internet\\
\rlab{1984--89} & Nodes pass one thousand, then one hundred thousand hosts &
Network effects begin\\
\rowcolor{tablealt}
\rlab{1989} & Tim Berners-Lee at CERN proposes cross-referencing documents
across the internet using hypertext & The concept of the World Wide Web\\
\rlab{6 Aug 1991} & The first web page goes online at CERN & The web becomes a
public medium\\
\rowcolor{tablealt}
\rlab{2000s on} & The dot-com boom, then Web 2.0 & User-generated content; the
global village\\
\end{mmlongtable}

\begin{exambox}[How to Frame This Answer]
If asked about the technology behind digital marketing, structure the answer in
three parts rather than listing dates: (i) the four-step adoption pattern of
Figure~\ref{fig:adoption}; (ii) the timeline of enabling technologies from
Table~\ref{tab:timeline}; and (iii) the present-day stack --- websites and
content management systems, search engines and their ranking algorithms, social
platforms and their feed algorithms, analytics and tag management, e-mail
service providers and CRM, mobile devices and apps, cloud infrastructure, and
increasingly machine learning for personalisation and pricing.
\end{exambox}

\section{Understanding the Digital Consumer}

\subsection{The myth and the reality}

There is a notion in marketing circles of mysterious, ethereal creatures who
exist in a hyper-connected cyber-world of their own --- enigmas who speak a
different language and are turning marketing on its head. The first thing to
realise about digital consumers is that there is basically no such thing. The
customers encountered online are the very same people who walk into the store,
telephone, or order from the catalogue.

They are doing exactly what people have done for thousands of years:
communicating with each other. That technology lets them do it faster, over
distance and across devices is perceived as dangerous and extraordinary, but
people have always talked to each other. What has genuinely changed is
\emph{consumer behaviour}, and it changed because of the pervasive, evocative
and enabling nature of digital technology.

\begin{keybox}[The Consumer Does Not Share Your Taxonomy]
Consumers do not care how marketers categorise their work. Above the line,
below the line, digital, traditional, experiential, linear, analogue, mobile,
direct, indirect --- all of these boxes are meaningless to them. All consumers
care about is the experience, and whether the marketing available to them helps
them make a more informed decision.
\end{keybox}

\subsection{Perceived anonymity}

\keytermas{Perceived anonymity}{perceived anonymity} is an online trait with a profound effect on
behaviour. It liberates consumers from the social shackles that bind them in
the physical world; online they act with scant regard for the social propriety
that holds sway in real life.

In a physical store, shoppers wait patiently for service and will endure a
less-than-perfect experience to get what they want. Online they will not. They
demand instant gratification and a flawless experience, and if you fail to
engage, retain and fulfil their expectations on demand they are gone --- the
only trace a solitary line in the server log.

\subsection{The five key traits of the online consumer}

\begin{mmlongtable}{The five traits of the online consumer}{tab:fivetraits}%
{@{}P{0.19\textwidth}P{0.38\textwidth}P{0.35\textwidth}@{}}
\rowcolor{ocre}\thd{Trait} & \thd{What it means} & \thd{What the marketer must therefore do}\\
\midrule
\endfirsthead
\rowcolor{ocre}\thd{Trait} & \thd{What it means} & \thd{What the marketer must therefore do}\\
\midrule
\endhead
\rlab{1. Comfortable with the medium} & Users have been online for years and
are increasingly web-savvy across all age groups; they use it efficiently and
do not hang around & Deliver what they want, fast; write scannable content\\
\rowcolor{tablealt}
\rlab{2. They want it all, and they want it now} & Accustomed to information on
demand from multiple sources simultaneously; their time is precious &
Design for scannability and instant gratification; front-load the value\\
\rlab{3. They are in control} & The web is no passive medium; Web 2.0 put users
firmly in charge & Make marketing user-centric, elective and permission-based,
with a real value proposition\\
\rowcolor{tablealt}
\rlab{4. They are fickle} & Transparency and immediacy erode brand loyalty;
competitors are one click and one search away & Build genuine trust; make sure
the total value proposition stacks up globally, not just the brand identity\\
\rlab{5. They are vocal} & Through reviews, blogs, social networks, forums and
communities they tell each other about good and bad experiences &
Harness the viral upside; monitor and respond, because the same mechanism
produces a backlash\\
\end{mmlongtable}

\subsection{Consumer 2.0 --- how Web 2.0 rewired behaviour}

\begin{itemize}
\item \sfb{Inter-connectivity} --- borderless peer-to-peer communication
      through social networks.
\item \sfb{On-demand} --- instant gratification, anytime and anywhere.
\item \sfb{Levelled information} --- unbiased research is freely available,
      shifting power decisively to the buyer.
\item \sfb{Rise of the \keytermas{prosumer}{prosumer}} --- customers now help
      shape and even dictate the products they consume.
\item \sfb{Relevance filtering} --- users block noise and customise their own
      feeds; irrelevant marketing is simply never seen.
\item \sfb{Niche aggregation} --- mass markets fragment into thousands of
      micro-communities.
\item \sfb{Micropublishing} --- anyone can publish opinions, reviews and blogs
      instantly.
\end{itemize}

\subsection{Online influencers}

\begin{definitionbox}[Influencer]
An \keyterm{influencer} is an early adopter and online opinion leader who,
through blogs, podcasts, forums and social networks, extols the virtues of
brands they like and denigrates those they find unsatisfactory.
\end{definitionbox}

Influencers matter because they hold the virtual ear of the online masses.
People read them, value their opinion and trust their judgement. Engage
positively and you recruit a team of powerful advocates; give them a negative
experience and the reverse happens, at the same scale. This is word-of-mouth
marketing on steroids.

\begin{examplebox}[The DoubleClick Test for Recognising an Influencer]
A person qualifies as an influencer if they strongly agree with three or more
of the following.
\begin{itemize}
\item They consider themselves expert in certain areas --- work, hobbies or
      interests.
\item People often ask their advice about purchases in areas where they are
      knowledgeable.
\item When they encounter a new product they like, they tend to recommend it to
      friends.
\item They have a large social circle and often refer people to one another
      based on interests.
\item They are active online, using blogs, social networking sites, e-mail,
      discussion groups and community boards to connect with peers.
\end{itemize}
\end{examplebox}

\begin{keybox}[Why Word of Mouth Dominates]
Roughly 90\% of people believe what their peers say about a brand, while only
about 48\% trust what businesses say directly. Americans value word-of-mouth
recommendations from friends and family about 41\% more than social media
recommendations. This asymmetry is the entire economic case for content
marketing, influencer marketing\ix{influencer marketing} and community building --- and the reason
earned media\ix{earned media} (\S\,3.2) is treated as the most credible of the three media
types.
\end{keybox}

\begin{summarybox}
Marketing management chooses target markets and gets, keeps and grows customers
by creating, delivering and communicating superior value; the mantra compresses
this to create, communicate, deliver, value, target, profit. Ten classes of
entity can be marketed. Needs are not created by marketers, wants are shaped,
and demand requires ability to pay. Value is defined by the audience, built
into the solution and judged in the result; performance above expectation
produces delight and therefore advocacy. A marketplace is physical, a market
space is digital, a metamarket spans industries. Firms have evolved through
production, product, selling and marketing concepts to holistic marketing.
Digital marketing is targeted, conversion-oriented, quantifiable and
interactive marketing through digital technologies --- not a separate
discipline, and not fundamentally about technology. Traditional marketing wins
on local reach and physical persistence and loses on measurement, cost and
interactivity; the transition to digital is sequenced, hybrid and
measurement-led. Push originates with the seller, pull with the buyer. Each of
the four Ps has a digital form. The consumer has not changed; consumer
behaviour has --- and the five traits, perceived anonymity and the influencer
are the practical consequences.
\end{summarybox}

%%%%%%%%%%%%%%%%%%%%%%%%%%%%%%%%%%%%%%%%%%%%%%%%%%%%%%%%%%%%%%%%%%%%%%%%%%%%%%
\chapter{Principles of Digital Marketing}
%%%%%%%%%%%%%%%%%%%%%%%%%%%%%%%%%%%%%%%%%%%%%%%%%%%%%%%%%%%%%%%%%%%%%%%%%%%%%%

\syllabusstrip{Principles of Digital Marketing.}

\begin{learningoutcomes}
\begin{itemize}
\item The five foundational pillars on which any digital strategy rests.
\item The internet paradox --- broadening reach and narrowing focus at once.
\item The seven governing principles: measurability, interactivity,
      personalisation, permission, integration, iteration, content centrality.
\item The six elements common to successful digital marketing campaigns.
\end{itemize}
\end{learningoutcomes}

\begin{keybox}[High Yield]
This chapter answers a recurring ten-mark question in full and a two-mark
question in one line.
\pyq{CIE-I, 1 Jul 2024 --- 10 M; SEE, Oct/Nov 2023 --- 2 M}
\end{keybox}

\section{The Five Foundational Pillars}

There is no one-size-fits-all digital strategy. Every business must bake its
own, and the best people to define it are the people who best know the
business. But every workable digital strategy rests on the same five
foundations, and each is a question rather than an assertion.

\begin{figure}[htbp]
\centering
\begin{tikzpicture}[font=\sffamily\footnotesize,
    pil/.style={draw=ocre,line width=1pt,fill=ocrepale,
                minimum width=2.02cm,minimum height=2.5cm,align=center,
                text width=1.75cm,font=\sffamily\scriptsize,text=slate,
                anchor=south},
    num/.style={circle,fill=ocre,text=white,inner sep=1.7pt,
                font=\sffamily\bfseries\scriptsize}]

  % base
  \fill[slate] (-0.35,-0.42) rectangle (11.05,-0.1);
  \node[anchor=north,text=slate,font=\sffamily\bfseries\scriptsize]
    at (5.35,-0.55) {A WORKABLE DIGITAL MARKETING STRATEGY};

  \node[pil] (p1) at (0.7,0)  {\textbf{Know your business}\\[2pt]
     Is it ready? Are the products suited? Is the skill set in place?};
  \node[pil] (p2) at (2.95,0) {\textbf{Know the competition}\\[2pt]
     What do they do right, wrong, and not at all?};
  \node[pil] (p3) at (5.2,0)  {\textbf{Know your customers}\\[2pt]
     Who are they, what do they want, how do they use technology?};
  \node[pil] (p4) at (7.45,0) {\textbf{Know what you want to achieve}\\[2pt]
     Clear, measurable, achievable goals};
  \node[pil] (p5) at (9.7,0)  {\textbf{Know how you're doing}\\[2pt]
     Measure, tweak, refine, re-measure};

  \foreach \n/\i in {p1/1,p2/2,p3/3,p4/4,p5/5}{
    \node[num] at ($(\n.north)+(0,0.22)$) {\i};}

  % capital
  \fill[ocre] (-0.35,2.62) rectangle (11.05,2.9);
\end{tikzpicture}
\caption[The five pillars of digital marketing strategy]%
        {The five pillars. Remove any one and the strategy rests on an
         assumption rather than on knowledge.}
\label{fig:fivepillars}
\end{figure}

\begin{enumerate}
\item \sfb{Know your business.} Is the business ready to embrace digital
      marketing? Are the products or services suited to online promotion? Is
      the right technology, skill set and infrastructure in place? How will
      digital marketing fit into existing business processes --- do those
      processes need to change, and are the staff ready for that change?
\item \sfb{Know the competition.} Who are the main competitors in the digital
      marketplace, and are they the same as the offline ones? What are they
      doing right, so it can be emulated; what are they doing wrong, so it can
      be learned from; and what are they not doing at all --- because that is
      the opportunity. Competition online can come from around the corner or
      around the globe, so never limit the analysis to local rivals.
\item \sfb{Know your customers.} Who are they and what do they want? Are you
      serving the same customer base online, or fishing in a new demographic?
      How do they use digital technology, and how can that be harnessed into an
      ongoing relationship?
\item \sfb{Know what you want to achieve.} If you don't know where you're
      going, you will never get there. Set clear, measurable, achievable goals
      --- online sales, targeted leads, brand awareness\ix{brand awareness} --- because goals are
      the yardsticks against which campaign progress is measured.
\item \sfb{Know how you're doing.} Everything online is trackable. Which
      channels deliver traffic? What are the conversion rates? How much of that
      traffic becomes tangible value? Measure, tweak, refine, re-measure.
\end{enumerate}

\section{The Internet Paradox}

\begin{keybox}[The Internet Paradox]
Unlike conventional mass media, the internet is unique in its capacity to
\textbf{both broaden the scope of marketing reach and narrow its focus at the
same time}. Digital channels let a firm transcend geography and time zones to
reach a far wider audience, while simultaneously honing the message with
laser-like precision onto very specific niche segments within that wider
market. Implemented well, it is an extraordinarily powerful combination.
\end{keybox}

\begin{figure}[htbp]
\centering
\begin{tikzpicture}[font=\sffamily\footnotesize]
  % broadening cone
  \fill[ocre!18] (0,0) -- (4.4,1.6) -- (4.4,-1.6) -- cycle;
  \draw[ocre,line width=0.9pt] (0,0) -- (4.4,1.6);
  \draw[ocre,line width=0.9pt] (0,0) -- (4.4,-1.6);
  \node[text=ocredark,font=\sffamily\bfseries\scriptsize,align=center]
    at (2.6,0) {BROADEN\\ THE REACH};
  \node[anchor=north,text=slate,font=\sffamily\scriptsize,align=center,
        text width=4.4cm] at (2.2,-1.75)
    {transcends geography and time zones\\ \textit{a far wider audience}};

  % narrowing cone
  \fill[deepblue!14] (10.6,0) -- (6.2,1.6) -- (6.2,-1.6) -- cycle;
  \draw[deepblue,line width=0.9pt] (10.6,0) -- (6.2,1.6);
  \draw[deepblue,line width=0.9pt] (10.6,0) -- (6.2,-1.6);
  \node[text=deepblue,font=\sffamily\bfseries\scriptsize,align=center]
    at (8.0,0) {NARROW\\ THE FOCUS};
  \node[anchor=north,text=slate,font=\sffamily\scriptsize,align=center,
        text width=4.4cm] at (8.4,-1.75)
    {age, location, interest, behaviour, intent\\ \textit{a very specific niche}};

  \node[circle,fill=greenok,text=white,inner sep=3pt,
        font=\sffamily\bfseries\scriptsize] at (5.3,0) {\hspace{-1pt}=\hspace{-1pt}};
  \node[anchor=south,text=greenok,font=\sffamily\bfseries\scriptsize]
    at (5.3,0.35) {SIMULTANEOUSLY};
\end{tikzpicture}
\caption[The internet paradox]%
        {The internet paradox. Conventional mass media forces a choice between
         reach and precision; digital channels do not.}
\label{fig:paradox}
\end{figure}

\section{The Governing Principles}

\begin{itemize}
\item \sfb{Measurability} --- every interaction leaves a quantifiable data
      trail. If it cannot be measured, it should not be a channel.
\item \sfb{Interactivity} --- marketing in the digital age is a dialogue, as
      much about listening as about telling.
\item \sfb{Personalisation} --- digital marketing allows uniquely tailored,
      ongoing relationships with individual customers rather than an
      undifferentiated mass.
\item \sfb{Permission and user control} --- the web is not a passive medium.
      Marketing must be user-centric, elective and permission-based, or the
      audience will actively disengage.
\item \sfb{Integration} --- channels must reinforce one another; a siloed
      campaign wastes the compounding effect.
\item \sfb{Iteration} --- digital marketing is never finished. It is a
      continuous engine.
\item \sfb{Content centrality} --- content is king. Produce powerful, genuine,
      compelling content and in time that content becomes the biggest promoter
      of the brand.
\end{itemize}

\begin{frameworkbox}[The Iteration Loop]
Execute strategy \ra{} measure analytics \ra{} tweak campaigns \ra{} refine
goals \ra{} execute strategy.
\end{frameworkbox}

\begin{exambox}[Two-Mark Answer]
``List any two principles of digital marketing.''
\pyq{SEE, Oct/Nov 2023 --- 2 M}
Any two of the five pillars --- know your business; know the competition; know
your customers; know what you want to achieve; know how you're doing --- or any
two of measurability, interactivity, personalisation and permission-based
targeting. Either set scores.
\end{exambox}

\section{Characteristics of a Successful Digital Marketing Campaign}

Campaigns differ enormously, but successful ones are built on the same few
elements.

\begin{enumerate}
\item \sfb{A quality company website.} What was a novelty a decade ago is a
      critical requirement today. A business website is a pseudo-B2B portal
      that lets businesses and customers around the world connect with you ---
      a company brochure available around the clock that projects the firm as a
      professional outfit.
\item \sfb{Social media presence.} Facebook, Instagram, LinkedIn, X, Pinterest
      and YouTube let a firm communicate across wide demographics of age, taste
      and culture. By actively promoting the company there, you get the world to
      talk about you.
\item \sfb{Blogging and forums.} Produce genuine, compelling content on the
      company blog and be a regular on relevant forums and discussion panels.
      Learn the art of connecting business themes into the content produced.
\item \sfb{Local listing services.} Register on Google Business Profile so the
      business appears in Google searches and on Google Maps; also list on Bing
      Places and equivalent directories. Joining is free and the reach
      translates into free promotion.
\item \sfb{E-mail marketing.} Done well and with good intentions, personalised
      e-mail speaks directly to the target audience, produces an immediate
      response and opens a line of communication with prospects --- fast, cheap
      promotion with measurable feedback.
\item \sfb{A dedicated digital marketing expert.} Hiring a professional lets
      management stay on top of promotional ideas while someone qualified
      implements the concepts that drive the expected results.
\end{enumerate}

\begin{summarybox}
Five pillars found any digital strategy: know your business, the competition,
your customers, your goals and your results. The internet paradox is that
digital channels broaden reach and narrow focus simultaneously, which
conventional mass media cannot do. Seven principles govern execution ---
measurability, interactivity, personalisation, permission, integration,
iteration and content centrality --- and iteration means the work is never
finished. Six elements recur in successful campaigns: a quality website, social
presence, blogging and forums, local listings, e-mail marketing and dedicated
expertise.
\end{summarybox}

%%%%%%%%%%%%%%%%%%%%%%%%%%%%%%%%%%%%%%%%%%%%%%%%%%%%%%%%%%%%%%%%%%%%%%%%%%%%%%
\chapter{Digital Marketing Channels}
%%%%%%%%%%%%%%%%%%%%%%%%%%%%%%%%%%%%%%%%%%%%%%%%%%%%%%%%%%%%%%%%%%%%%%%%%%%%%%

\syllabusstrip{Digital Marketing Channels.}

\begin{learningoutcomes}
\begin{itemize}
\item The three senses in which marketing uses the word ``channel''.
\item The paid, owned and earned media taxonomy, and what each is for.
\item The campaign ecosystem as a hub with spokes, and the job of each spoke.
\item Affiliate marketing --- definition, payment models, advantages and risks.
\item How to choose between SEO, PPC and social media on a limited budget.
\end{itemize}
\end{learningoutcomes}

\section{The Three Senses of ``Channel''}

Before listing platforms, note that \emph{channel} in marketing has three
distinct senses, and a question about channels may mean any of them.

\begin{itemize}
\item \sfb{Communication channels} deliver messages to, and receive messages
      from, target buyers.
\item \sfb{Distribution channels} display, sell or deliver the physical product
      or service to the buyer.
\item \sfb{Service channels} carry out transactions with prospective buyers ---
      warehouses, transporters, banks, payment gateways, insurers.
\end{itemize}

\section{Paid, Owned and Earned Media}

This is the single most useful way to classify digital channels, and it is
examined directly.
\pyq{CIE-I, 16 Apr 2026 --- 10 M}

\begin{figure}[htbp]
\centering
\begin{tikzpicture}[font=\sffamily\footnotesize]
  \begin{scope}[opacity=0.42,transparency group]
    \fill[ocre]     (-1.35,0.55) circle (2.15);
    \fill[deepblue] ( 1.35,0.55) circle (2.15);
    \fill[greenok]  ( 0.00,-1.6) circle (2.15);
  \end{scope}
  \node[text=ocredark,font=\sffamily\bfseries\small]  at (-2.5,1.55) {OWNED};
  \node[text=deepblue,font=\sffamily\bfseries\small]  at ( 2.5,1.55) {PAID};
  \node[text=greenok,font=\sffamily\bfseries\small]   at ( 0.0,-3.3) {EARNED};

  \node[text=slate,font=\sffamily\scriptsize,align=center,text width=2.3cm]
    at (-2.35,0.45) {website, blog,\\ app, e-mail list,\\ brand profiles};
  \node[text=slate,font=\sffamily\scriptsize,align=center,text width=2.3cm]
    at ( 2.35,0.45) {display ads, PPC,\\ social ads,\\ sponsorships};
  \node[text=slate,font=\sffamily\scriptsize,align=center,text width=2.4cm]
    at ( 0.0,-2.15) {word of mouth, reviews,\\ press, backlinks};

  \node[text=white,font=\sffamily\bfseries\scriptsize,align=center]
    at (0,-0.35) {the\\ campaign};

  % legend of roles
  \node[anchor=west,align=left,font=\sffamily\scriptsize,text=slate]
    at (4.0,0.6)
    {\textcolor{ocredark}{\textbf{Owned} --- the foundation.}\\
     No rent, full control, compounds.\\[5pt]
     \textcolor{deepblue}{\textbf{Paid} --- the accelerator.}\\
     Instant reach; stops when\\ spending stops.\\[5pt]
     \textcolor{greenok}{\textbf{Earned} --- the validator.}\\
     Most credible, least\\ controllable.};
\end{tikzpicture}
\caption[The paid, owned and earned media taxonomy]%
        {Every digital channel is paid, owned or earned. A campaign that uses
         only one of the three is either renting all its reach, talking only to
         people it already has, or hoping.}
\label{fig:poe}
\end{figure}

\begin{mmlongtable}{Paid, owned and earned media}{tab:poe}%
{@{}P{0.108\textwidth}P{0.275\textwidth}P{0.313\textwidth}P{0.216\textwidth}@{}}
\rowcolor{ocre}\thd{Media type} & \thd{Definition} & \thd{Examples} & \thd{Role}\\
\midrule
\endfirsthead
\rowcolor{ocre}\thd{Media type} & \thd{Definition} & \thd{Examples} & \thd{Role}\\
\midrule
\endhead
\rlab{Owned} & Assets entirely controlled by the brand & Company website, blog,
mobile app, e-mail list, brand social profiles, brochures & The foundation ---
no rent, full control, compounds over time\\
\rowcolor{tablealt}
\rlab{Paid} & Capital deployed for guaranteed reach & Display advertising, paid
search and PPC, social media advertising, sponsorships, magazine and television
advertising, influencer promotions & The accelerator --- instant reach, stops
the moment spending stops\\
\rlab{Earned} & Audience-driven reach that was not paid for & Word of mouth,
buzz, viral sharing, press coverage, reviews, influencer mentions, backlinks &
The validator --- most credible, least controllable\\
\end{mmlongtable}

\begin{keybox}[Rented Against Owned --- The Apple Tree Rule]
Advertising is a rented treadmill. Content is an owned asset. The moment
payment for advertising stops, traffic falls to zero; content keeps working.

\smallskip
\textbf{The apple tree rule:} it takes time to plant the seed and grow the tree,
but once mature that single tree bears fruit for years without anyone having to
buy the apples.
\end{keybox}

\section{The Campaign Ecosystem --- Hub and Spokes}

A coherent digital campaign is not a list of disconnected activities. It is a
hub-and-spoke system whose hub is a quality company website --- the round-the-clock
portal and the ultimate conversion destination. Every other channel exists to
feed it.

\begin{figure}[htbp]
\centering
\begin{tikzpicture}[font=\sffamily\footnotesize,
    spoke/.style={draw=ocre!70,line width=0.7pt,fill=ocrepale,
                  rounded corners=3pt,align=center,inner sep=3.5pt,
                  font=\sffamily\scriptsize,text=slate,text width=1.95cm}]

  \node[draw=slate,line width=1.2pt,fill=white,circle,minimum size=2.6cm,
        align=center,font=\sffamily\bfseries\scriptsize,text=slate] (hub) at (0,0)
        {THE HUB\\[1pt] company\\ website};

  \foreach \ang/\name/\lbl in {%
      90/s1/{SEO},
      54/s2/{PPC / SEM},
      18/s3/{Social media},
     -18/s4/{Content, blogs, forums},
     -54/s5/{E-mail},
     -90/s6/{Affiliate},
    -126/s7/{Display and remarketing},
    -162/s8/{Mobile},
     162/s9/{Video},
     126/s10/{Local listings}}{
     \node[spoke] (\name) at (\ang:4.05cm) {\lbl};
     \draw[-{Latex[length=1.9mm]},ocre!70,line width=0.7pt] (\name) -- (hub);}
\end{tikzpicture}
\caption[The campaign ecosystem]%
        {The hub-and-spoke campaign ecosystem. Every spoke has one job: to
         deliver qualified attention to the hub, where conversion happens.}
\label{fig:hubspoke}
\end{figure}

\begin{mmlongtable}{The job of each spoke}{tab:spokes}%
{@{}P{0.30\textwidth}P{0.63\textwidth}@{}}
\rowcolor{ocre}\thd{Spoke} & \thd{Primary job}\\
\midrule
\endfirsthead
\rowcolor{ocre}\thd{Spoke} & \thd{Primary job}\\
\midrule
\endhead
\rlab{Search Engine Optimisation} & Capture high-intent organic demand and
drive it to the hub\\
\rowcolor{tablealt}
\rlab{Search Engine Marketing / PPC} & Buy immediate visibility on high-value
commercial queries\\
\rlab{Social Media Marketing} & Manage global chatter, build awareness, and
engineer community\\
\rowcolor{tablealt}
\rlab{Content Marketing --- blogging and forums} & Produce the compelling,
genuine material that every other channel distributes\\
\rlab{E-mail Marketing} & Own the relationship; deliver personalised,
rapid-response communication\\
\rowcolor{tablealt}
\rlab{Affiliate Marketing} & Pay partners only for results --- traffic, leads or
sales actually delivered\\
\rlab{Display and Network Advertising} & Build reach and remarket to prior
visitors\\
\rowcolor{tablealt}
\rlab{Mobile Marketing} & Reach the customer on the device that is always with
them\\
\rlab{Video Marketing} & Build trust and dwell time; dual search-and-social
distribution\\
\rowcolor{tablealt}
\rlab{Local Listings} & Be discoverable in ``near me'' and map searches\\
\rlab{Influencer Marketing} & Borrow trust that the brand has not yet earned\\
\end{mmlongtable}

\section{Affiliate Marketing}

\begin{definitionbox}[Affiliate Marketing]
\keytermas{Affiliate marketing}{affiliate marketing} is a type of performance-based marketing in which
a business rewards affiliates --- partners or individuals --- for driving
traffic, leads or sales to the company's products or services through the
affiliate's own marketing efforts.
\pyq{CIE-I, 3 Apr 2025 --- 2 M}
\end{definitionbox}

\begin{examplebox}[Amazon Associates]
A blogger places a tracked Amazon link inside a product review and earns a
commission on every purchase made through it. Amazon pays nothing for the
review itself, nothing for the traffic that does not convert, and a fixed
percentage on the traffic that does.
\end{examplebox}

\begin{itemize}
\item \sfb{Payment models.} CPA --- cost per acquisition or sale; CPL --- cost
      per lead; CPC --- cost per click.
\item \sfb{Advantages.} Risk is transferred to the affiliate, because payment
      follows results; reach extends into audiences the brand does not own; and
      costs are variable rather than fixed.
\item \sfb{Risks.} Brand-safety exposure if affiliates use aggressive or
      misleading tactics; potential for fraud; and margin erosion if
      commissions are set too high.
\end{itemize}

\section{Choosing Between Channels}

\begin{exambox}[The Scenario Question]
``Your company has a limited marketing budget and needs to choose between
Search Engine Optimization, Pay Per Click, and Social Media Marketing to
promote a new product. Based on the product being a high-tech gadget aimed at
young professionals, which channel(s) would you prioritize and why?''
\pyq{CIE-I, 1 Jul 2024 --- 10 M}

\smallskip
A strong answer compares the three on a fixed set of criteria and \emph{then
commits} to a recommendation. A comparison without a recommendation loses marks.
\end{exambox}

\begin{mmlongtable}{SEO, PPC and social media marketing compared}{tab:channelchoice}%
{@{}P{0.16\textwidth}P{0.25\textwidth}P{0.25\textwidth}P{0.25\textwidth}@{}}
\rowcolor{ocre}\thd{Criterion} & \thd{SEO} & \thd{PPC} & \thd{Social Media Marketing}\\
\midrule
\endfirsthead
\rowcolor{ocre}\thd{Criterion} & \thd{SEO} & \thd{PPC} & \thd{Social Media Marketing}\\
\midrule
\endhead
\rlab{Speed of results} & Slow --- months & Immediate & Fast for reach, slow
for trust\\
\rowcolor{tablealt}
\rlab{Cost model} & Free per click; costs time and resources & Charged per
click; stops when the budget stops & Organic free; paid boosts optional\\
\rlab{Longevity} & Infinite --- compounds & Zero --- ends with the campaign &
Short, hours to days per post\\
\rowcolor{tablealt}
\rlab{Buyer intent} & High --- the user is actively searching & High ---
keyword-triggered & Low to medium --- interruption-based\\
\rlab{Predictability} & High once established & High & Low --- relies on viral
spikes\\
\rowcolor{tablealt}
\rlab{Best suited to} & Long-term compounding demand & Launches, promotions,
high-margin conversions & Awareness, community, visual products\\
\end{mmlongtable}

\begin{exambox}[Model Recommendation]
For a high-tech gadget aimed at young professionals on a limited budget:

\begin{itemize}
\item \sfb{about 50\% to PPC and paid social} for the launch window. The
      product is new, so nobody is searching for it by name yet --- organic
      search cannot deliver demand that does not exist.
\item \sfb{about 30\% to social media marketing}, weighted to Instagram,
      YouTube and LinkedIn, because young professionals discover technology
      through video reviews and peer validation.
\item \sfb{about 20\% to SEO and content} --- comparison articles, specification
      pages, ``best gadget for X'' content --- as the compounding asset that
      reduces paid dependence over the following two quarters.
\end{itemize}

\smallskip
\textbf{Justify with the intent argument:} PPC and social \emph{create} the
demand; SEO \emph{harvests} it once it exists. That single sentence is what
distinguishes a first-class answer from a list.
\end{exambox}

\begin{summarybox}
Channel means communication, distribution or service. Every digital channel is
paid, owned or earned: owned is the foundation, paid the accelerator, earned the
validator. A campaign is a hub of one quality website fed by spokes, each with a
single job. Affiliate marketing pays only for delivered results, on CPA, CPL or
CPC terms, transferring risk at the cost of brand-safety exposure. Channel
choice is decided by speed, cost model, longevity, buyer intent\ix{buyer intent} and
predictability --- and a scenario answer must end in a weighted recommendation
justified by buyer intent.
\end{summarybox}

%%%%%%%%%%%%%%%%%%%%%%%%%%%%%%%%%%%%%%%%%%%%%%%%%%%%%%%%%%%%%%%%%%%%%%%%%%%%%%
\chapter{Tools to Create a Buyer Persona}
%%%%%%%%%%%%%%%%%%%%%%%%%%%%%%%%%%%%%%%%%%%%%%%%%%%%%%%%%%%%%%%%%%%%%%%%%%%%%%

\syllabusstrip{Tools to Create Buyer Persona.}

\begin{learningoutcomes}
\begin{itemize}
\item What a buyer persona is, and how many a business needs.
\item The five-step process for building one, in the order the scheme expects.
\item The seven factor categories to capture, with an illustration throughout.
\item The persona builder canvas --- four quadrants of questions.
\item A complete worked persona.
\item The toolset, and the division of labour between hard data and listening.
\end{itemize}
\end{learningoutcomes}

\begin{keybox}[The Most Repeated Ten-Mark Question in This Course]
The adventure-tourism travel agency scenario is set verbatim in two separate
papers. Learn the process, the factors, and be ready to invent a named
example persona on the spot.
\pyq{CIE-I, 3 Apr 2025 --- 10 M; CIE-I, 16 Apr 2026 --- 10 M}
\end{keybox}

\section{What a Buyer Persona Is}

\begin{definitionbox}[Buyer Persona]
A \keyterm{buyer persona} is a semi-fictional, research-based profile of the
ideal customer --- given a name, a face, a backstory and a set of goals, pain
points and behaviours --- built from real data about existing customers rather
than from assumption.
\end{definitionbox}

The governing rule is that you cannot hit a target you have not defined. The
working ideal is four to six distinct personas per business --- for a B2B
purchase, for example, an Executive Sponsor, a Decision Maker and an End User,
who have different goals and raise different objections about the same product.

\begin{warnbox}[Semi-Fictional, Not Fictional]
The word ``semi-fictional'' does real work. The \emph{name and the photograph}
are invented so the profile is memorable and shareable. The \emph{goals, pain
points, objections and behaviours} are not invented --- they come from surveys,
analytics, CRM records and review mining. A persona built entirely from
imagination is a stereotype, and it will mislead every decision downstream.
\end{warnbox}

\section{The Five-Step Process}

\begin{frameworkbox}[The Official Five Steps]
1. Research your current customers \ra{} 2. Segment your audience \ra{}
3. Gather data \ra{} 4. Identify goals and pain points \ra{} 5. Build the
persona --- then validate and update it continuously.
\end{frameworkbox}

\begin{enumerate}
\item \sfb{Research your current customers.} Analyse the best existing clients
      --- those who buy often, leave reviews and engage with content. Conduct
      surveys or interviews, online or through post-purchase feedback forms.
\item \sfb{Segment your audience.} Divide customers into categories based on
      shared traits. For a travel agency: by travel habits, destination
      preferences and trip budgets.
\item \sfb{Gather data.} Use tools such as Google Analytics\ix{Google Analytics} and social media
      insights to identify patterns in who actually interacts with the website
      and the social pages --- as distinct from who management imagines does.
\item \sfb{Identify goals and pain points.} For an adventure-tourism customer,
      the desire for adventure set directly against safety concerns. Personas
      are most useful where the goal and the pain point are in tension, because
      that tension is what the marketing has to resolve.
\item \sfb{Build the persona.} Give it a name, a photograph and a backstory to
      make it relatable --- ``Adventurous Ananya'', ``Explorer Ethan'' --- and
      attach the key data points. Then validate and update it using continuing
      customer feedback and trend data.
\end{enumerate}

\section{Factors to Capture}

\begin{mmlongtable}{Persona factors, with an adventure-tourism illustration}{tab:personafactors}%
{@{}P{0.16\textwidth}P{0.36\textwidth}P{0.38\textwidth}@{}}
\rowcolor{ocre}\thd{Category} & \thd{Factors to capture} & \thd{Adventure-tourism illustration}\\
\midrule
\endfirsthead
\rowcolor{ocre}\thd{Category} & \thd{Factors to capture} & \thd{Adventure-tourism illustration}\\
\midrule
\endhead
\rlab{Demographics} & Age, gender, income level, occupation, education,
location, family status & 20--40 years; working professionals and students;
income sufficient to afford travel; metro-based\\
\rowcolor{tablealt}
\rlab{Psychographics} & Values, attitudes, hobbies, aspirations, risk appetite &
Interest in trekking, scuba diving, skydiving; fitness and outdoor lifestyle;
risk-taking attitude; prefers unique experiences over luxury\\
\rlab{Online behaviour} & Platforms used, search patterns, content preference,
device, booking behaviour & Instagram, YouTube and travel blogs; searches
``best trekking places''; prefers videos, reels and reviews; books online
packages via mobile apps\\
\rowcolor{tablealt}
\rlab{Category behaviour} & Frequency, budget tier, group against solo,
preferred variants & Travels two or three times a year; budget against premium;
solo or small group; mountains, beaches, forests\\
\rlab{Goals} & What success looks like for them & Unique, thrilling, shareable
experiences\\
\rowcolor{tablealt}
\rlab{Pain points} & What blocks the purchase & Safety concerns, budget
constraints, lack of reliable information, distrust of agencies\\
\rlab{Objections} & What they say when they say no & ``Is it safe?'' ``Is the
guide certified?'' ``What if I have to cancel?''\\
\end{mmlongtable}

\section{The Persona Builder Canvas}

\begin{figure}[htbp]
\centering
\begin{tikzpicture}[font=\sffamily\footnotesize,
    q/.style={draw=ocre,line width=0.8pt,minimum width=5.6cm,
              minimum height=2.55cm,align=left,text width=5.1cm,
              font=\sffamily\scriptsize,text=slate,inner sep=7pt}]

  \node[q,fill=ocrepale] (q1) at (0,0)      {\textbf{1. BACKGROUND \& GOALS}\\[3pt]
     What are their business goals and challenges?\\ What is their team
     structure and job role?};
  \node[q,fill=white]    (q2) at (5.75,0)   {\textbf{2. INFORMATION SOURCES}\\[3pt]
     Where do they consume content?\\ Which social channels do they actively
     use --- and which do they avoid?};
  \node[q,fill=white]    (q3) at (0,-2.62)  {\textbf{3. MESSAGING \& TOPICS}\\[3pt]
     Which marketing messages speak to them?\\ What content topics capture
     their interest?};
  \node[q,fill=ocrepale] (q4) at (5.75,-2.62){\textbf{4. PURCHASE ROLE \& OBJECTIONS}\\[3pt]
     How much influence do they have in the decision?\\ What objections do they
     raise during the sales process?};
\end{tikzpicture}
\caption[The persona builder canvas]%
        {The persona builder canvas. Quadrants 1 and 3 establish who they are
         and what moves them; quadrants 2 and 4 establish where to reach them
         and what will stop the sale.}
\label{fig:personacanvas}
\end{figure}

\section{A Worked Persona}

\begin{examplebox}[Persona: ``Adventure Arun'']
\begin{tabular}{@{}>{\sffamily\bfseries\scriptsize\color{slate}}P{0.16\textwidth}P{0.72\textwidth}@{}}
Name        & Adventure Arun\\
Age         & 28\\
Occupation  & IT professional, Bengaluru\\
Interests   & Trekking, camping, weekend rides\\
Behaviour   & Follows travel influencers on Instagram and YouTube; researches on
              travel blogs; books entirely online, usually on mobile, usually
              three to four weeks ahead\\
Goal        & Unique, thrilling experiences he can share\\
Pain points & Safety and affordability; not knowing whether the operator is
              credible\\
What wins him & Certified-guide credentials, real customer video testimonials,
              transparent inclusions, easy cancellation, instalment options\\
\end{tabular}
\end{examplebox}

Note what the last row does. ``What wins him'' converts the persona from a
description into an instruction: each item is a specific asset the marketing
team must now produce. A persona that does not end in such a list has not
finished its job.

\section{The Toolset}

\begin{mmlongtable}{Tools for building a buyer persona}{tab:personatools}%
{@{}P{0.18\textwidth}P{0.36\textwidth}P{0.36\textwidth}@{}}
\rowcolor{ocre}\thd{Category} & \thd{Tools} & \thd{What they give you}\\
\midrule
\endfirsthead
\rowcolor{ocre}\thd{Category} & \thd{Tools} & \thd{What they give you}\\
\midrule
\endhead
\rlab{Analytics} & Google Analytics 4, Google Search Console & Demographics,
interests, devices, locations, acquisition channels, on-site behaviour\\
\rowcolor{tablealt}
\rlab{Social insights} & Meta Audience Insights, Instagram Insights, LinkedIn
Page Analytics, X Analytics, YouTube Studio & Follower demographics, active
hours, content resonance\\
\rlab{Survey and voice of customer} & Google Forms, Typeform, SurveyMonkey,
Hotjar polls, post-purchase feedback forms & Stated goals, objections and
satisfaction\\
\rowcolor{tablealt}
\rlab{Persona builders} & HubSpot Make My Persona, Xtensio, Userforge, Semrush
Persona & A structured template and a shareable persona document\\
\rlab{Behavioural and qualitative mining} & Quora, Reddit, Amazon reviews,
YouTube comments, competitor review sections & The exact behavioural problems
and questions the audience is trying to solve, in their own words\\
\rowcolor{tablealt}
\rlab{CRM and sales data} & HubSpot, Zoho, Salesforce, sales-call notes & Real
objections, deal-cycle length, decision roles\\
\end{mmlongtable}

\begin{keybox}[Hard Data Tells You Who; Social Listening Tells You Why]
Use Google Analytics and Meta Audience Insights for hard data on age, location
and income. Use Quora, Reddit and Amazon reviews to extract the behavioural
material --- pain points, common questions, motivations, desired outcomes,
objections and influences --- that the audience is actively trying to resolve.
Neither source substitutes for the other: analytics cannot tell you a motive,
and a forum thread cannot tell you a market size.
\end{keybox}

\begin{summarybox}
A buyer persona is a semi-fictional, research-based profile of the ideal
customer; four to six per business is the working ideal. Build it in five steps
--- research current customers, segment, gather data, identify goals and pain
points, then construct and continuously validate. Capture seven categories of
factor: demographics, psychographics, online behaviour, category behaviour,
goals, pain points and objections. The canvas organises the questions into
background and goals, information sources, messaging and topics, and purchase
role and objections. Analytics and social insights supply who; surveys, review
mining and CRM records supply why. A finished persona ends in a list of assets
the team must produce.
\end{summarybox}

%%%%%%%%%%%%%%%%%%%%%%%%%%%%%%%%%%%%%%%%%%%%%%%%%%%%%%%%%%%%%%%%%%%%%%%%%%%%%%
\chapter{Competitor Research Tools}
%%%%%%%%%%%%%%%%%%%%%%%%%%%%%%%%%%%%%%%%%%%%%%%%%%%%%%%%%%%%%%%%%%%%%%%%%%%%%%

\syllabusstrip{Competitor Research Tools.}

\begin{learningoutcomes}
\begin{itemize}
\item The seven things to look for in a competitor analysis.
\item The toolset, and what each tool is specifically good at.
\item Why the content gap\ix{content gap} is the most valuable output of the exercise.
\end{itemize}
\end{learningoutcomes}

``Know the competition'' is the second pillar of digital strategy (\S\,2.1).
Competition online can come from around the corner or around the globe, so the
analysis must never be limited to local rivals.

\section{What to Look For}

\begin{figure}[htbp]
\centering
\begin{tikzpicture}[font=\sffamily\footnotesize,
    it/.style={draw=ocre!65,line width=0.7pt,fill=ocrepale,rounded corners=3pt,
               align=center,text width=2.55cm,inner sep=4pt,
               font=\sffamily\scriptsize,text=slate}]
  \node[draw=slate,line width=1.1pt,fill=white,circle,minimum size=2.35cm,
        align=center,font=\sffamily\bfseries\scriptsize,text=slate] (c) at (0,0)
        {COMPETITOR\\ ANALYSIS};
  \foreach \ang/\n/\lbl in {%
      103/a/{Organic and paid keywords},
      52/b/{Backlink profile},
      0/c2/{\textbf{Content gaps}},
     -52/d/{Traffic sources and volume},
    -103/e/{Social share of voice},
     -155/f/{Ad creatives, offers, landing pages},
     155/g/{Pricing and positioning}}{
     \node[it] (\n) at (\ang:3.9cm) {\lbl};
     \draw[ocre!65,line width=0.7pt] (\n) -- (c);}
  % highlight the content-gap node
  \node[draw=greenok,line width=1.1pt,fill=greenpale,rounded corners=3pt,
        align=center,text width=2.55cm,inner sep=4pt,
        font=\sffamily\bfseries\scriptsize,text=greenok] at (0:3.9cm)
        {Content gaps\\[1pt]{\mdseries\scriptsize where opportunity lives}};
\end{tikzpicture}
\caption[What to look for in a competitor analysis]%
        {The seven dimensions of competitor analysis. Six describe what rivals
         are doing; the seventh --- the content gap --- describes what none of
         them is doing, and that is where the opportunity is.}
\label{fig:competitor}
\end{figure}

\begin{itemize}
\item Which keywords they rank for organically, and which they pay for.
\item Their \sfb{backlink profile} --- who links to them, and why.
\item Their \sfb{content gaps} --- topics with demonstrated demand that nobody
      has covered well. This is where opportunity lives.
\item Their traffic sources and volume, and how these are trending.
\item Their \sfb{social share of voice} --- conversations about them against
      conversations about you.
\item Their advertising creatives, offers and landing pages.
\item Their pricing, positioning and value proposition\ix{value proposition}.
\end{itemize}

\section{The Toolset}

\begin{mmlongtable}{Competitor research tools}{tab:competitortools}%
{@{}P{0.27\textwidth}P{0.66\textwidth}@{}}
\rowcolor{ocre}\thd{Tool} & \thd{Primary use}\\
\midrule
\endfirsthead
\rowcolor{ocre}\thd{Tool} & \thd{Primary use}\\
\midrule
\endhead
\rlab{SEMrush} & All-round competitive intelligence --- organic and paid
keywords, traffic estimates, advertising copy, backlink gap, Keyword Magic Tool\\
\rowcolor{tablealt}
\rlab{Ahrefs} & Backlink analysis on the deepest link index; Content Gap and Top
Pages reports\\
\rlab{Moz / Link Explorer} & Domain Authority and Page Authority scoring, link
metrics\\
\rowcolor{tablealt}
\rlab{SimilarWeb} & Traffic estimates, traffic sources, audience overlap,
referral analysis\\
\rlab{SpyFu} & Historical competitor PPC keyword and advertising-spend data\\
\rowcolor{tablealt}
\rlab{BuzzSumo} & Which content in a topic earned the most shares and backlinks
--- content-gap discovery\\
\rlab{Google Trends} & Relative search interest over time, seasonality,
regional variation, rising queries\\
\rowcolor{tablealt}
\rlab{Google Keyword Planner} & Search volume and competition data straight from
the advertising platform\\
\rlab{Meta Ad Library, Google Ads Transparency Center} & The competitor's live
advertising creatives, free of charge\\
\rowcolor{tablealt}
\rlab{Social listening --- Hootsuite, Brandwatch, Mention, Sprout Social} &
Share of voice, sentiment, brand mention tracking\\
\end{mmlongtable}

\section{Turning Findings into Decisions}

A competitor report that ends in observations has failed. Each finding should
resolve into one of three decisions.

\begin{itemize}
\item \sfb{Emulate} --- they are doing something that works, and it is
      transferable. Copy the mechanism, not the creative.
\item \sfb{Avoid} --- they are doing something that visibly is not working.
      Their sunk cost is your free experiment.
\item \sfb{Occupy} --- nobody is doing it, and there is demand. This is the
      content gap, and it is the only one of the three that produces a durable
      advantage, because the other two put you where competitors already are.
\end{itemize}

\begin{summarybox}
Competitor analysis examines organic and paid keywords, backlinks, content gaps,
traffic sources and trends, social share of voice\ix{share of voice}, advertising creatives and
landing pages, and pricing and positioning. SEMrush\ix{SEMrush} and Ahrefs carry the
keyword and backlink work, SimilarWeb the traffic estimates, SpyFu the
historical advertising data, BuzzSumo the content performance, Google Trends\ix{Google Trends} the
seasonality, and the ad transparency libraries the live creatives at no cost.
Every finding must resolve into emulate, avoid or occupy --- and only occupying
a genuine content gap produces durable advantage.
\end{summarybox}

%%%%%%%%%%%%%%%%%%%%%%%%%%%%%%%%%%%%%%%%%%%%%%%%%%%%%%%%%%%%%%%%%%%%%%%%%%%%%%
\chapter{Website Analysis Tools}
%%%%%%%%%%%%%%%%%%%%%%%%%%%%%%%%%%%%%%%%%%%%%%%%%%%%%%%%%%%%%%%%%%%%%%%%%%%%%%

\syllabusstrip{Website Analysis Tools.}

\begin{learningoutcomes}
\begin{itemize}
\item The two halves of website analysis: technical health and behavioural
      insight.
\item The toolset by category, and what each tool actually measures.
\item The six mechanisms by which website analysis improves customer experience
      and conversion.
\end{itemize}
\end{learningoutcomes}

\begin{keybox}[High Yield]
This chapter answers a compulsory Question~2 sub-division in the Semester End
Examination.
\pyq{SEE, Jun/Jul 2025, Q2a --- 8 M}
\end{keybox}

\section{The Two Halves}

Website analysis has two halves, and a complete answer names both. \sfb{Technical
health} asks whether the site works --- whether it loads fast, whether it can be
crawled and indexed, whether links resolve. \sfb{Behavioural insight} asks what
people do once it works --- where they hesitate, where they scroll, where they
abandon.

\begin{figure}[htbp]
\centering
\begin{tikzpicture}[font=\sffamily\footnotesize,
    half/.style={draw=ocre,line width=1pt,rounded corners=4pt,
                 minimum width=5.35cm,minimum height=1.05cm,
                 align=center,font=\sffamily\bfseries\small,text=slate},
    lst/.style={align=left,text width=5.0cm,font=\sffamily\scriptsize,
                text=slate,anchor=north}]
  \node[half,fill=ocrepale] (t) at (0,0)     {TECHNICAL HEALTH};
  \node[half,fill=deepbluepale,draw=deepblue] (b) at (5.7,0) {BEHAVIOURAL INSIGHT};
  \node[lst] at (0,-0.72) {%
     \textit{does the site work?}\\[3pt]
     load speed and Core Web Vitals\\
     crawlability and indexing\\
     broken links and redirect chains\\
     duplicate titles and missing meta\\
     code validity and accessibility};
  \node[lst] at (5.7,-0.72) {%
     \textit{what do people do on it?}\\[3pt]
     acquisition channels and sessions\\
     engagement and funnel drop-off\\
     heatmaps, scroll maps, click maps\\
     session recordings and rage clicks\\
     A/B and multivariate test results};
  \node[draw=slate,line width=0.9pt,fill=white,rounded corners=3pt,
        minimum width=11.05cm,align=center,font=\sffamily\bfseries\scriptsize,
        text=slate,anchor=north] at (2.85,-3.35)
     {a diagnosis, not a dashboard: \textnormal{\itshape one named
      friction point and one change to test}};
\end{tikzpicture}
\caption[The two halves of website analysis]%
        {Technical health and behavioural insight answer different questions.
         Reporting either alone produces numbers without a decision.}
\label{fig:websiteanalysis}
\end{figure}

\section{The Toolset}

\begin{mmlongtable}{Website analysis tools}{tab:websitetools}%
{@{}P{0.29\textwidth}P{0.64\textwidth}@{}}
\rowcolor{ocre}\thd{Tool} & \thd{What it analyses}\\
\midrule
\endfirsthead
\rowcolor{ocre}\thd{Tool} & \thd{What it analyses}\\
\midrule
\endhead
\rlab{Google Analytics 4} & Sessions, users, acquisition channels, engagement,
conversions, funnels, audience demographics\\
\rowcolor{tablealt}
\rlab{Google Search Console} & Impressions, clicks, average position, indexing
status, Core Web Vitals, crawl errors\\
\rlab{Google PageSpeed Insights / Lighthouse} & Load speed, Core Web Vitals
(LCP, INP, CLS), accessibility, SEO and best-practice audits\\
\rowcolor{tablealt}
\rlab{GTmetrix / Pingdom} & Page load waterfall, asset weight, server response
time\\
\rlab{Hotjar / Microsoft Clarity / Crazy Egg} & Heatmaps, scroll maps, click
maps and session recordings --- what users actually do\\
\rowcolor{tablealt}
\rlab{Screaming Frog SEO Spider} & Full technical crawl --- broken links,
duplicate titles, redirect chains, missing meta\\
\rlab{SEMrush Site Audit / Ahrefs Site Audit} & Automated technical SEO health
score and prioritised issue list\\
\rowcolor{tablealt}
\rlab{VWO / Optimizely} & A/B testing and multivariate testing of pages and
elements\\
\rlab{Google Tag Manager} & Central deployment and management of all tracking
tags without touching code\\
\rowcolor{tablealt}
\rlab{W3C Validator / WAVE} & Code validity and accessibility conformance\\
\end{mmlongtable}

\section{How Website Analysis Improves Experience and Conversion}

\begin{exambox}[Model Answer --- 8 Marks]
\pyq{SEE, Jun/Jul 2025, Q2a --- 8 M}
Name the tools by category, then give the six mechanisms. The mechanisms, not
the tool names, are what earn the marks.

\begin{enumerate}
\item \sfb{Diagnoses friction.} Heatmaps and session recordings expose where
      users hesitate, rage-click or abandon --- behaviour no aggregate metric
      would reveal.
\item \sfb{Fixes speed.} Slow pages are the single largest cause of
      abandonment, so speed audits convert directly into retained sessions.
\item \sfb{Reveals the drop-off point in the funnel.} Funnel reports show which
      step loses people, so the fix is targeted rather than cosmetic.
\item \sfb{Validates changes objectively} through A/B testing\ix{A/B testing} instead of
      opinion or seniority.
\item \sfb{Segments experience by device and channel}, exposing faults that
      averages hide --- a mobile-only checkout bug, for instance, is invisible
      in a blended conversion rate.
\item \sfb{Aligns content with intent} by showing which pages attract traffic
      that never converts --- the classic ``traffic without intent'' problem.
\end{enumerate}
\end{exambox}

\begin{summarybox}
Website analysis splits into technical health --- speed, Core Web Vitals\ix{Core Web Vitals},
crawlability, broken links, code validity --- and behavioural insight ---
acquisition, engagement, funnels, heatmaps, recordings and split tests. Google
Analytics 4 and Search Console are the base layer; PageSpeed Insights\ix{PageSpeed Insights} and
GTmetrix cover performance; Hotjar and Clarity cover behaviour; Screaming Frog\ix{Screaming Frog}
and the site-audit tools cover crawl\ix{crawl} health; VWO and Optimizely run the
experiments; Tag Manager deploys the tracking. Analysis improves conversion by
diagnosing friction, fixing speed, locating funnel drop-off, validating changes
by test, segmenting by device and channel, and aligning content with intent.
\end{summarybox}

%%%%%%%%%%%%%%%%%%%%%%%%%%%%%%%%%%%%%%%%%%%%%%%%%%%%%%%%%%%%%%%%%%%%%%%%%%%%%%
\chapter{Content Marketing --- Concepts and Strategies}
%%%%%%%%%%%%%%%%%%%%%%%%%%%%%%%%%%%%%%%%%%%%%%%%%%%%%%%%%%%%%%%%%%%%%%%%%%%%%%

\syllabusstrip{Content Marketing: Content Marketing Concepts \& Strategies.}

\begin{learningoutcomes}
\begin{itemize}
\item What content marketing is, and the one rule that governs it.
\item A century of evidence that the strategy predates the internet.
\item The compounding traffic curve, and why it beats the advertising sawtooth.
\item The five traits that distinguish elite content marketers.
\item The seven-point quality gate for content --- and the second accepted list.
\end{itemize}
\end{learningoutcomes}

\section{Definition and the Core Idea}

\begin{definitionbox}[Content Marketing]
\keytermas{Content marketing}{content marketing} is a strategic marketing approach focused on
creating and distributing valuable, relevant and consistent content to attract
and retain a clearly defined audience --- and, ultimately, to drive profitable
customer action.
\end{definitionbox}

Four words in that definition are load-bearing. \emph{Valuable} means the
audience would want it even if the brand were removed from it.
\emph{Relevant} means it connects to what they are doing now.
\emph{Consistent} means published to a rhythm, because a single asset does not
compound. And \emph{profitable customer action} means the objective is
commercial --- content marketing is not publishing for its own sake.

\begin{keybox}[The Governing Principle]
\textbf{Don't sell the product; solve the problem.} Great content integrates
the product into the customer's natural pursuit of their own interests, rather
than interrupting that pursuit to talk about the product.
\end{keybox}

\section{A Century of Proof}

Content marketing is routinely presented as a consequence of the internet. It
is not. The three canonical examples all predate it by decades.

\begin{mmlongtable}{Content marketing before the internet}{tab:centuryproof}%
{@{}P{0.09\textwidth}P{0.30\textwidth}P{0.55\textwidth}@{}}
\rowcolor{ocre}\thd{Year} & \thd{Example} & \thd{Strategy}\\
\midrule
\endfirsthead
\rowcolor{ocre}\thd{Year} & \thd{Example} & \thd{Strategy}\\
\midrule
\endhead
\rlab{1895} & John Deere's \textit{The Furrow} & Educate the farmer --- a
magazine about farming, not about tractors\\
\rowcolor{tablealt}
\rlab{1900} & The Michelin Guide & Inspire the driver to travel farther for a
three-star meal, creating the need for more tyres\\
\rlab{1930s} & Soap operas & Entertain the homemaker with radio dramas sponsored
by soap companies\\
\end{mmlongtable}

\begin{keybox}[The Synthesis]
The medium changes. The strategy does not. Michelin did not advertise tyres; it
made driving farther desirable. John Deere did not advertise ploughs; it made
farmers better at farming. That is still exactly what content marketing does.
\end{keybox}

\section{The Compounding Traffic Curve}

Paid advertising produces a sawtooth: spend, spike, stop, collapse. Content
produces a curve that starts slower and then rises continuously, because each
asset keeps earning traffic long after it is published, and new assets stack on
top of old ones.

\begin{figure}[htbp]
\centering
\begin{tikzpicture}
\begin{axis}[
  width=12.2cm, height=6.1cm,
  font=\sffamily\footnotesize,
  xlabel={months since launch}, ylabel={monthly traffic},
  xlabel style={font=\sffamily\scriptsize,text=slate},
  ylabel style={font=\sffamily\scriptsize,text=slate},
  axis lines=left,
  axis line style={rulegrey},
  tick label style={font=\sffamily\scriptsize,text=slatelight},
  xmin=0, xmax=18, ymin=0, ymax=115,
  ytick=\empty, xtick={0,3,6,9,12,15,18},
  clip=false,
  legend style={font=\sffamily\scriptsize,draw=rulegrey,fill=white,
                at={(0.02,0.97)},anchor=north west,text=slate}]

  % paid advertising: sawtooth that collapses whenever spend pauses
  \addplot[deepblue,thick,const plot] coordinates {
    (0,0) (1,42) (2,44) (3,43) (4,4) (5,3)
    (6,45) (7,46) (8,44) (9,4) (10,3)
    (11,44) (12,46) (13,45) (14,4) (15,3) (16,44) (17,45) (18,44)};
  \addlegendentry{paid advertising --- rented}

  % content: compounding
  \addplot[ocre,very thick,smooth,domain=0:18,samples=90] {0.335*x^2};
  \addlegendentry{content --- owned}

  \node[anchor=west,text=deepblue,font=\sffamily\scriptsize,align=left]
    at (axis cs:4.15,17) {spend pauses,\\ traffic collapses};
  \node[anchor=east,text=ocredark,font=\sffamily\scriptsize,align=right]
    at (axis cs:17.6,103) {each asset keeps\\ earning; new assets\\ stack on old};
\end{axis}
\end{tikzpicture}
\caption[The compounding content curve against the advertising sawtooth]%
        {Paid advertising buys a plateau that ends with the budget. Content buys
         a curve. Content is slower to the first result and the only one of the
         two that accumulates.}
\label{fig:compounding}
\end{figure}

\begin{keybox}[The Apple Tree Rule]
It takes time to plant the seed and grow the tree. But once mature, that single
tree continually bears fruit for years, without requiring anyone to buy the
apples. Advertising is a rented treadmill; content is an owned asset.
\end{keybox}

\section{The Five Traits of Elite Content Marketers}

\begin{enumerate}
\item \sfb{Documented strategy} --- they operate from a written plan, not from
      intuition.
\item \sfb{Dedicated leadership} --- a specific named person is accountable for
      the content engine.
\item \sfb{Expanded tactics} --- they use a wider variety of content formats
      than their peers.
\item \sfb{Multi-platform distribution} --- they leverage more channels than
      their peers.
\item \sfb{Budget allocation} --- they commit a higher percentage of the total
      marketing budget to content.
\end{enumerate}

\section{Characteristics of Good Content --- The Seven-Point Quality Gate}

\begin{keybox}[Guaranteed Short Question]
Two slightly different lists are accepted; learn the first and know the second.
\pyq{CIE-I, 3 Apr 2025; CIE-I, 16 Apr 2026; SEE, Jun/Jul 2025 --- 2 M each}
\end{keybox}

\begin{mmlongtable}{The seven-point content quality gate}{tab:qualitygate}%
{@{}P{0.20\textwidth}P{0.73\textwidth}@{}}
\rowcolor{ocre}\thd{Trait} & \thd{The diagnostic question}\\
\midrule
\endfirsthead
\rowcolor{ocre}\thd{Trait} & \thd{The diagnostic question}\\
\midrule
\endhead
\rlab{Credible}     & Will they believe it?\\
\rowcolor{tablealt}
\rlab{Shareable}    & Will the audience want to share it?\\
\rlab{Useful or fun} & Does it provide genuine utility or entertainment?\\
\rowcolor{tablealt}
\rlab{Interesting}  & Does it break through apathy? --- necessarily subjective\\
\rlab{Relevant}     & Does it connect to their current context?\\
\rowcolor{tablealt}
\rlab{Different}    & Does it stand out from the noise?\\
\rlab{On-brand}     & Is it authentically aligned with our identity?\\
\end{mmlongtable}

\begin{exambox}[The Second Accepted List]
The 2026 scheme of valuation accepted any four of: \textbf{relevant, accurate,
engaging, clear, consistent, valuable}. If the question says ``mention any
four'', either list scores full marks --- but the seven-point gate above is the
one taught in the lecture notes, and it is the safer answer because each trait
carries a diagnostic question you can expand if asked to explain.
\end{exambox}

\begin{summarybox}
Content marketing creates and distributes valuable, relevant, consistent content
to attract and retain a defined audience and drive profitable action; the rule
is solve the problem, don't sell the product. The strategy predates the internet
--- John Deere in 1895, Michelin in 1900, sponsored radio drama in the 1930s.
Advertising produces a sawtooth that collapses when spending stops; content
produces a compounding curve, which is the apple tree rule. Elite practitioners
share five traits: a documented strategy, dedicated leadership, expanded
tactics, multi-platform distribution and higher budget share. Good content is
credible, shareable, useful or fun, interesting, relevant, different and
on-brand.
\end{summarybox}

%%%%%%%%%%%%%%%%%%%%%%%%%%%%%%%%%%%%%%%%%%%%%%%%%%%%%%%%%%%%%%%%%%%%%%%%%%%%%%
\chapter{Planning, Creating, Distributing and Promoting Content}
%%%%%%%%%%%%%%%%%%%%%%%%%%%%%%%%%%%%%%%%%%%%%%%%%%%%%%%%%%%%%%%%%%%%%%%%%%%%%%

\syllabusstrip{Planning, Creating, Distributing \& Promoting Content.}

\begin{learningoutcomes}
\begin{itemize}
\item The eight-step content marketing process, in three phases.
\item Each step in detail: goals, audience mapping, ideation, creation,
      distribution, amplification, evaluation, improvement.
\item How to map content to funnel stage and to the 5A customer path.
\item The complete content flywheel\ix{content flywheel}.
\end{itemize}
\end{learningoutcomes}

\begin{keybox}[The Framework That Answers Every Scenario Question]
Every ``design a content marketing strategy for company X'' question in this
course is answered by walking through these eight steps in order. Memorise the
step names.
\pyq{CIE-I, 3 Apr 2025 --- 10 M; CIE-I, 1 Jul 2024 --- 10 M; CIE-I, 16 Apr 2026 --- 10 M}
\end{keybox}

\section{The Eight Steps in Three Phases}

\begin{frameworkbox}[The Official Scheme Answer]
1. Goal Setting \ra{} 2. Audience Mapping \ra{} 3. Content Ideation and
Planning \ra{} 4. Content Creation \ra{} 5. Content Distribution \ra{}
6. Content Amplification \ra{} 7. Content Marketing Evaluation \ra{}
8. Content Marketing Improvement --- and then back to step 1.

\smallskip
It is a cycle, not a checklist.
\end{frameworkbox}

\begin{figure}[htbp]
\centering
\begin{tikzpicture}[font=\sffamily\footnotesize,
    st/.style={draw=ocre,line width=0.8pt,fill=white,rounded corners=3pt,
               minimum height=0.82cm,align=center,text width=2.35cm,
               font=\sffamily\scriptsize,text=slate},
    ph/.style={rounded corners=4pt,inner sep=8pt,align=center,
               font=\sffamily\bfseries\scriptsize}]

  % ---- phase bands
  \fill[ocrepale,rounded corners=5pt]      (-0.05,-0.62) rectangle (7.75,0.62);
  \fill[deepbluepale,rounded corners=5pt]  (-0.05,-2.62) rectangle (7.75,-1.38);
  \fill[greenpale,rounded corners=5pt]     (-0.05,-4.62) rectangle (5.2,-3.38);

  \node[anchor=west,text=ocredark,font=\sffamily\bfseries\scriptsize]
     at (8.05,0)     {PHASE 1 --- ARCHITECTURE};
  \node[anchor=west,text=deepblue,font=\sffamily\bfseries\scriptsize]
     at (8.05,-2)    {PHASE 2 --- MECHANICS};
  \node[anchor=west,text=greenok,font=\sffamily\bfseries\scriptsize]
     at (8.05,-4)    {PHASE 3 --- TELEMETRY};
  \node[anchor=west,text=slatelight,font=\sffamily\scriptsize\itshape]
     at (8.05,-0.42) {why, for whom, and what};
  \node[anchor=west,text=slatelight,font=\sffamily\scriptsize\itshape]
     at (8.05,-2.42) {build it, ship it, spread it};
  \node[anchor=west,text=slatelight,font=\sffamily\scriptsize\itshape]
     at (8.05,-4.42) {did it work, and what next};

  % ---- steps
  \node[st] (s1) at (1.2,0)    {\textbf{1}\ \ Goal setting};
  \node[st] (s2) at (3.85,0)   {\textbf{2}\ \ Audience mapping};
  \node[st] (s3) at (6.5,0)    {\textbf{3}\ \ Ideation \& planning};
  \node[st] (s4) at (1.2,-2)   {\textbf{4}\ \ Creation};
  \node[st] (s5) at (3.85,-2)  {\textbf{5}\ \ Distribution};
  \node[st] (s6) at (6.5,-2)   {\textbf{6}\ \ Amplification};
  \node[st] (s7) at (1.2,-4)   {\textbf{7}\ \ Evaluation};
  \node[st] (s8) at (3.85,-4)  {\textbf{8}\ \ Improvement};

  \foreach \a/\b in {s1/s2,s2/s3,s4/s5,s5/s6,s7/s8}{
     \draw[-{Latex[length=2mm]},ocre,line width=0.85pt] (\a)--(\b);}
  \draw[-{Latex[length=2mm]},ocre,line width=0.85pt]
     (s3.south) -- ++(0,-0.42) -| (s4.north);
  \draw[-{Latex[length=2mm]},ocre,line width=0.85pt]
     (s6.south) -- ++(0,-0.42) -| (s7.north);

  % ---- the return loop
  \draw[-{Latex[length=2mm]},ocredark,line width=0.85pt,dashed]
     (s8.east) -- ++(0.9,0) -- ++(0,4.9) -- (s1.north);
  \node[anchor=south,text=ocredark,font=\sffamily\scriptsize\itshape]
     at (2.5,0.95) {the cycle restarts with revised goals};
\end{tikzpicture}
\caption[The eight-step content marketing process]%
        {The eight steps in three phases. Steps 7 and 8 are not a report at the
         end --- they are the input to the next cycle's step 1.}
\label{fig:eightsteps}
\end{figure}

\begin{mmlongtable}{The three phases}{tab:threephases}%
{@{}P{0.24\textwidth}P{0.36\textwidth}P{0.33\textwidth}@{}}
\rowcolor{ocre}\thd{Phase} & \thd{Steps} & \thd{Purpose}\\
\midrule
\endfirsthead
\rowcolor{ocre}\thd{Phase} & \thd{Steps} & \thd{Purpose}\\
\midrule
\endhead
\rlab{Phase 1 --- Architecture} & Goal setting, audience mapping, ideation and
planning & Decide why, for whom, and what\\
\rowcolor{tablealt}
\rlab{Phase 2 --- Mechanics} & Creation, distribution, amplification &
Actually build it, and get it in front of people\\
\rlab{Phase 3 --- Telemetry} & Evaluation, improvement & Find out whether it
worked, and make the next cycle better\\
\end{mmlongtable}

\section{Step 1 --- Goal Setting}

The primary objective always splits into two branches, and every subsequent
decision inherits from this split.

\begin{itemize}
\item \sfb{Brand-building objective} --- awareness, recall, reputation,
      authority, share of voice.
\item \sfb{Sales-growth objective} --- leads, conversions, revenue, repeat
      purchase.
\end{itemize}

Three reflective questions found the goal: what content will be valuable to the
company; how can the content tell a story about the brand; and how will the
content marketing strategy actually be executed.

\begin{warnbox}[The First Step Is Not ``Create Content'']
Asked ``what is the first step of content marketing?'', the official answer is
\textbf{define goals and identify the target audience}. Answering ``create
content'' scores nothing.
\pyq{CIE-I, 16 Apr 2026 --- 2 M}
\end{warnbox}

\section{Step 2 --- Audience Mapping}

Audience mapping is built outward from internal psychology, not inward from
demographics, and the order matters.

\begin{figure}[htbp]
\centering
\begin{tikzpicture}[font=\sffamily\footnotesize]
  \fill[ocre!12] (0,0) circle (3.5);
  \fill[ocre!25] (0,0) circle (2.35);
  \fill[ocre!45] (0,0) circle (1.2);
  \draw[ocre,line width=0.8pt] (0,0) circle (3.5);
  \draw[ocre,line width=0.8pt] (0,0) circle (2.35);
  \draw[ocre,line width=0.8pt] (0,0) circle (1.2);

  \node[text=slate,font=\sffamily\bfseries\scriptsize,align=center] at (0,0)
     {1. ANXIETIES\\ \& DESIRES};
  \node[text=slate,font=\sffamily\bfseries\scriptsize,align=center] at (0,1.78)
     {2. PAIN POINTS};
  \node[text=slate,font=\sffamily\bfseries\scriptsize,align=center] at (0,2.95)
     {3. PROFILE \& PERSONA};

  \node[anchor=west,text=slatelight,font=\sffamily\scriptsize,align=left,
        text width=4.9cm] at (4.0,1.3)
     {\textbf{The order is the point.}\\[3pt]
      The psychological foundation sits at the core; the pain points are how it
      manifests as real problems; the persona is only the visible artefact
      built on top.};
  \node[anchor=west,text=warnred,font=\sffamily\scriptsize,align=left,
        text width=4.9cm] at (4.0,-1.35)
     {\textbf{The common failure.}\\[3pt]
      Most failed content strategies invert this --- starting from a demographic
      and never reaching the anxiety.};
\end{tikzpicture}
\caption[The three layers of audience mapping]%
        {Audience mapping works from the inside out. Chapter~4 gives the full
         persona method that produces the outermost layer.}
\label{fig:audiencemapping}
\end{figure}

\section{Step 3 --- Content Ideation and Planning}

Ideation must be constrained by a disciplined roadmap, or it degenerates into
random posting. Four layers produce the plan.

\begin{mmlongtable}{The four layers of the content plan}{tab:fourlayers}%
{@{}P{0.20\textwidth}P{0.73\textwidth}@{}}
\rowcolor{ocre}\thd{Layer} & \thd{Output}\\
\midrule
\endfirsthead
\rowcolor{ocre}\thd{Layer} & \thd{Output}\\
\midrule
\endhead
\rlab{Strategic layer}  & The content theme --- the territory the brand will own\\
\rowcolor{tablealt}
\rlab{Structural layer} & The content formats and mix --- which media, in what
ratio\\
\rlab{Temporal layer}   & The storyline and calendar --- what publishes when, and
in what narrative order\\
\rowcolor{tablealt}
\rlab{Result}           & \textbf{The content road map}\\
\end{mmlongtable}

\section{Step 4 --- Content Creation}

The pivot question is who creates the content, and when. Two resource branches
feed a single production schedule: in-house creators and external agencies.
In-house wins on brand authenticity and subject knowledge; agencies win on
capacity, production quality and speed.

\subsection{Aligning format to funnel stage}

\begin{mmlongtable}{Content format by funnel stage}{tab:formatfunnel}%
{@{}P{0.20\textwidth}P{0.44\textwidth}P{0.29\textwidth}@{}}
\rowcolor{ocre}\thd{Stage} & \thd{Formats} & \thd{KPI}\\
\midrule
\endfirsthead
\rowcolor{ocre}\thd{Stage} & \thd{Formats} & \thd{KPI}\\
\midrule
\endhead
\rlab{Awareness} & Infographics, social videos, podcasts & Reach and shares\\
\rowcolor{tablealt}
\rlab{Engagement} & How-to blogs, video tutorials & Dwell time and video views\\
\rlab{Evaluation} & Case studies, e-books, whitepapers & Leads and downloads\\
\rowcolor{tablealt}
\rlab{Purchase and retention} & Newsletters, native advertising, product updates
& Conversion and open rates\\
\end{mmlongtable}

\subsection{Anatomy of a top-ranking page}

\begin{itemize}
\item \sfb{H1 title} --- click-worthy, with the exact-match keyword.
\item \sfb{Introduction} --- comprehensive, answering the user's core intent
      immediately rather than after four paragraphs of preamble.
\item \sfb{Structure} --- H2 and H3 subheadings containing long-tail keyword\ix{long-tail keyword}
      variants.
\item \sfb{Formatting} --- short, scannable paragraphs and bulleted lists.
\item \sfb{Rich media} --- embedded videos and infographics, to increase dwell
      time.
\item \sfb{Architecture} --- internal links to keep users in your ecosystem,
      plus credible outbound links to build trust.
\end{itemize}

\section{Step 5 --- Content Distribution}

Distribution runs on the owned, paid and earned trinity of \S\,3.2. Choose the
mix using the distribution matrix below, which trades predictability against
lifespan.

\begin{mmlongtable}{The distribution ecosystem matrix}{tab:distribution}%
{@{}P{0.17\textwidth}P{0.17\textwidth}P{0.17\textwidth}P{0.40\textwidth}@{}}
\rowcolor{ocre}\thd{Channel} & \thd{Predictability} & \thd{Lifespan} & \thd{Core characteristics}\\
\midrule
\endfirsthead
\rowcolor{ocre}\thd{Channel} & \thd{Predictability} & \thd{Lifespan} & \thd{Core characteristics}\\
\midrule
\endhead
\rlab{Google SEO} & High & Infinite & High buyer intent, search-driven\\
\rowcolor{tablealt}
\rlab{YouTube} & Medium & Long & Dual search-and-social algorithm; massive
visual trust\\
\rlab{Social media} & Low & Short --- hours & Flatlining traffic curve; relies
on viral spikes\\
\rowcolor{tablealt}
\rlab{E-mail} & High & Medium & Direct ownership; highest conversion and
retention\\
\end{mmlongtable}

\section{Step 6 --- Content Amplification}

Amplification creates compounding momentum by activating network effects. The
central question is how to leverage content assets and interact with customers.
Three nodes answer it.

\begin{enumerate}
\item \sfb{Creating communities} --- owned groups, forums and user communities
      that redistribute content for you.
\item \sfb{Use of buzzers} --- employees, partners and enthusiasts who seed and
      share content early, before any algorithm has decided whether it deserves
      reach.
\item \sfb{Use of influencers} --- borrowing the trust of established opinion
      leaders (\S\,1.9.5).
\end{enumerate}

\section{Mapping Content to the 5A Customer Path}

\begin{figure}[htbp]
\centering
\begin{tikzpicture}[font=\sffamily\footnotesize,
    a/.style={draw=ocre,line width=0.9pt,fill=ocrepale,
              signal,signal to=east,signal pointer angle=115,
              minimum height=1.05cm,text width=1.55cm,align=center,
              font=\sffamily\bfseries\scriptsize,text=slate,inner sep=3pt}]
  \foreach \i/\lbl/\req/\x in {%
      1/AWARE/{visible}/0,
      2/APPEAL/{relatable}/2.25,
      3/ASK/{searchable}/4.5,
      4/ACT/{actionable}/6.75,
      5/ADVOCATE/{shareable}/9.0}{
    \node[a] at (\x,0) {\lbl};
    \node[anchor=north,text=ocredark,font=\sffamily\scriptsize\itshape]
      at (\x+0.15,-0.68) {\req};}
  \node[anchor=north,text=slate,font=\sffamily\scriptsize,align=center,
        text width=12cm] at (4.6,-1.35)
    {found where attention already is \ \textbullet\ \ feels like it is about
     them \ \textbullet\ \ answers what they type\\
     makes the next step obvious \ \textbullet\ \ worth passing on};
\end{tikzpicture}
\caption[Content requirements along the 5A customer path]%
        {The 5A path. Each stage imposes a different requirement on the content,
         and content that satisfies one stage will not satisfy the next.}
\label{fig:fivea}
\end{figure}

\begin{examplebox}[A Content-Driven Customer Journey]
\begin{tabular}{@{}>{\sffamily\bfseries\scriptsize\color{slate}}P{0.19\textwidth}P{0.33\textwidth}P{0.36\textwidth}@{}}
\textcolor{ocre}{\thd{}} & \multicolumn{1}{@{}l}{\sffamily\bfseries\scriptsize\color{slate}What the customer thinks}
 & \multicolumn{1}{@{}l}{\sffamily\bfseries\scriptsize\color{slate}What they consume}\\
\midrule
Problem aware & ``The kids need something to climb on.'' & Searches ``best home
climbing equipment''\\
Solution research & ``Mini jungle gyms exist.'' & Reads a ``top ten jungle
gyms'' blog post\\
Product evaluation & ``Will my oldest fit?'' & Watches a video product review\\
Purchase & ``This model works.'' & Reads marketplace reviews and buys\\
\end{tabular}

\smallskip
\textbf{The architect's note:} the brand must be intentionally present at every
stepping stone. If your content is not there, a competitor's is.
\end{examplebox}

\section{Steps 7 and 8 --- Evaluation and Improvement}

The evaluation question is how successful the content marketing campaign was,
and it is answered through two measurement protocols: tracking specific
content-marketing metrics --- KPIs such as traffic, engagement rate, lead
generation and conversion tied to individual assets --- and determining overall
objective achievement, the holistic impact on brand awareness, reputation and
revenue, which loops back to the brand-building and sales-growth goals set in
step 1.

Improvement then pulls on exactly two levers.

\begin{itemize}
\item \sfb{Lever 1 --- content theme change and improvement}, which adjusts the
      architecture of Phase 1.
\item \sfb{Lever 2 --- distribution and amplification improvement}, which
      adjusts the mechanics of Phase 2.
\end{itemize}

Two practices keep the loop running: \sfb{A/B testing}, continuously split-testing
headline variants, call-to-action tone, image choices and content length; and
\sfb{historical audits}, regularly updating older high-ranking content with
fresh data, new internal links and current examples, to protect traffic that
already exists.

\begin{figure}[htbp]
\centering
\begin{tikzpicture}[font=\sffamily\footnotesize]
  \def\R{3.15}
  \foreach \i/\lbl/\ang in {%
      1/{Audience\\ research}/90,
      2/{Keyword\\ strategy}/18,
      3/{Content\\ creation}/-54,
      4/{Multi-channel\\ distribution}/-126,
      5/{Meaningful\\ measurement}/162}{
    \node[draw=ocre,line width=0.9pt,fill=ocrepale,rounded corners=4pt,
          align=center,text width=2.1cm,inner sep=5pt,
          font=\sffamily\bfseries\scriptsize,text=slate] (n\i) at (\ang:\R) {\lbl};}
  \foreach \a/\b in {1/2,2/3,3/4,4/5,5/1}{
    \draw[-{Latex[length=2.3mm]},ocre,line width=1pt]
      (n\a) to[bend left=17] (n\b);}
  \node[circle,draw=greenok,line width=1pt,fill=greenpale,
        align=center,text width=1.85cm,font=\sffamily\bfseries\scriptsize,
        text=greenok,inner sep=2pt] at (0,0) {compounding revenue growth};
\end{tikzpicture}
\caption[The content flywheel]%
        {The content flywheel. Each phase powers the next, and measurement
         optimises the following cycle rather than merely reporting on the last
         one. Content is never finished; it is an asset requiring continuous
         maintenance.}
\label{fig:flywheel}
\end{figure}

\begin{exambox}[Worked Scenario --- ``EcoChef'' Eco-Friendly Kitchen Appliances]
\pyq{CIE-I, 1 Jul 2024 --- 10 M}
Apply the eight-step framework explicitly.

\begin{itemize}
\item \sfb{Goals.} Build awareness of EcoChef among urban eco-conscious
      households; generate 5,000 e-mail subscribers in six months; achieve a
      3\% landing-page conversion rate.
\item \sfb{Personas.} Sustainability-minded young families; first-home buyers.
\item \sfb{Themes.} Sustainable cooking, energy savings, kitchen-to-table
      living, product demonstrations.
\item \sfb{Content types.} SEO blogs (``how much energy does an induction hob
      save''), recipe videos and reels, comparison infographics, a downloadable
      energy-savings calculator, webinars with chefs, customer-generated content.
\item \sfb{Distribution.} Owned --- website, blog, e-mail. Paid --- Google Ads on
      high-intent appliance keywords, Instagram and video advertising. Earned ---
      sustainability influencers, press, product reviews.
\item \sfb{Engagement.} A community group, a recipe contest, responsive social
      customer care, a nurture e-mail sequence.
\item \sfb{Measurement.} Reach, engagement rate, leads, conversions and
      content-attributed revenue --- feeding back into theme and distribution
      improvement.
\end{itemize}
\end{exambox}

\begin{exambox}[Worked Scenario --- Organic Health and Wellness Startup]
\pyq{CIE-I, 16 Apr 2026 --- 10 M}
The official scheme structures the answer in five headed parts.

\begin{enumerate}
\item \sfb{Objectives} --- build brand awareness; educate customers on organic
      lifestyle; generate leads and conversions.
\item \sfb{Content themes} --- organic living and its benefits; health and
      nutrition; sustainability and eco-friendly practices; product awareness;
      customer testimonials and expert advice.
\item \sfb{Content formats} --- blogs and articles for SEO; videos and reels
      (workouts, product demonstrations); infographics (diet charts); podcasts
      (wellness talks); e-books and webinars; user-generated content\ix{user-generated content} and
      reviews.
\item \sfb{Distribution platforms} --- \emph{owned}: website, blog, e-mail;
      \emph{earned}: social shares, influencer mentions, public relations;
      \emph{paid}: social advertising, Google Ads, influencer promotions.
\item \sfb{Metrics} --- awareness: reach, impressions, website traffic;
      engagement: likes, shares, comments, watch time; leads: sign-ups,
      downloads; conversions: sales and conversion rate; retention: repeat
      customers and loyalty.
\end{enumerate}
\end{exambox}

\begin{summarybox}
Content marketing runs as an eight-step cycle in three phases: architecture
(goals, audience mapping, ideation), mechanics (creation, distribution,
amplification) and telemetry (evaluation, improvement). Goals split into
brand-building and sales-growth, and everything downstream inherits from that
split. Audience mapping works inward-out from anxieties to pain points to
persona. The plan has strategic, structural and temporal layers producing a road
map. Format follows funnel stage. Distribution trades predictability against
lifespan across SEO, video, social and e-mail. Amplification uses communities,
buzzers and influencers. The 5A path demands content that is visible, relatable,
searchable, actionable and shareable in turn. Evaluation feeds two improvement
levers --- theme and distribution --- and the flywheel restarts.
\end{summarybox}

%%%%%%%%%%%%%%%%%%%%%%%%%%%%%%%%%%%%%%%%%%%%%%%%%%%%%%%%%%%%%%%%%%%%%%%%%%%%%%
\chapter{Optimizing Website UX and Landing Pages}
%%%%%%%%%%%%%%%%%%%%%%%%%%%%%%%%%%%%%%%%%%%%%%%%%%%%%%%%%%%%%%%%%%%%%%%%%%%%%%

\syllabusstrip{Optimize Website UX \& Landing Pages.}

\begin{learningoutcomes}
\begin{itemize}
\item Krug's first law of usability\ix{usability}.
\item The seven usability attributes and their test questions.
\item How people actually use the web --- scanning, satisficing, muddling
      through.
\item Billboard design: designing for scanning rather than reading.
\item The eight essentials of a converting landing page.
\end{itemize}
\end{learningoutcomes}

Traffic that arrives at a badly designed page is wasted money. This material is
shared ground with Unit~IV (web design) and Unit~V (conversion optimization);
what follows is the content-marketing view of it.

\section{Krug's First Law of Usability}

\begin{keybox}[Don't Make Me Think]
A person of average --- or below average --- ability must be able to work out how
to accomplish a task on the page. It must be self-evident and self-explanatory.
The good news is that usability is not rocket science.
\end{keybox}

\section{The Seven Usability Attributes}

\begin{mmlongtable}{The seven usability attributes}{tab:usability}%
{@{}P{0.18\textwidth}P{0.75\textwidth}@{}}
\rowcolor{ocre}\thd{Attribute} & \thd{Test question}\\
\midrule
\endfirsthead
\rowcolor{ocre}\thd{Attribute} & \thd{Test question}\\
\midrule
\endhead
\rlab{Useful}     & Does it do what people actually need?\\
\rowcolor{tablealt}
\rlab{Learnable}  & Can they figure it out?\\
\rlab{Memorable}  & Do they have to relearn it every visit?\\
\rowcolor{tablealt}
\rlab{Effective}  & Does it get the job done?\\
\rlab{Efficient}  & Does it do so in a reasonable amount of time?\\
\rowcolor{tablealt}
\rlab{Desirable}  & Do people want it?\\
\rlab{Delightful} & Is it enjoyable to use?\\
\end{mmlongtable}

\section{How People Actually Use the Web}

There is a gap between how designers imagine people use websites --- carefully
reading instructions like a novel --- and how they actually do: \keyterm{scanning},
\keyterm{satisficing} and muddling through. Satisficing means choosing the first
option that plausibly works rather than the best one available.

\begin{figure}[htbp]
\centering
\begin{tikzpicture}[font=\sffamily\footnotesize,
    pg/.style={draw=rulegrey,line width=0.8pt,fill=white,
               minimum width=4.1cm,minimum height=4.5cm}]
  % --- what we think
  \node[pg] (a) at (0,0) {};
  \node[anchor=north,text=slate,font=\sffamily\bfseries\scriptsize]
     at (0,2.85) {WHAT WE THINK THEY DO};
  \foreach \y in {1.85,1.6,1.35,1.1,0.85,0.6,0.35,0.1,-0.15,-0.4,-0.65,-0.9,-1.15,-1.4,-1.65,-1.9}{
     \draw[slatelight,line width=1.1pt] (-1.85,\y) -- (1.85,\y);}
  \node[anchor=north,text=slatelight,font=\sffamily\scriptsize\itshape]
     at (0,-2.45) {read every word, in order};

  % --- what they do
  \node[pg] (b) at (5.6,0) {};
  \node[anchor=north,text=slate,font=\sffamily\bfseries\scriptsize]
     at (5.6,2.85) {WHAT THEY ACTUALLY DO};
  \draw[ocre,line width=1.6pt] (3.75,1.85) -- (7.45,1.85);
  \draw[ocre,line width=1.6pt] (3.75,1.5) -- (6.6,1.5);
  \draw[ocre,line width=1.6pt] (3.75,0.6) -- (7.0,0.6);
  \draw[ocre,line width=1.6pt] (3.75,-0.5) -- (5.9,-0.5);
  \draw[ocre,line width=1.6pt] (3.75,-1.55) -- (6.9,-1.55);
  \draw[-{Latex[length=2mm]},deepblue,line width=1pt,dashed]
     (3.9,1.7) -- (6.9,1.35) -- (4.1,0.45) -- (6.6,-0.62) -- (4.2,-1.7);
  \node[anchor=north,text=ocredark,font=\sffamily\scriptsize\itshape]
     at (5.6,-2.45) {scan for what looks clickable};
\end{tikzpicture}
\caption[Reading against scanning]%
        {Users do not read pages; they scan them for the first thing that looks
         like it will get the job done. Pages must be designed for the right-hand
         behaviour.}
\label{fig:scanning}
\end{figure}

Why don't people read instructions? Three reasons.

\begin{enumerate}
\item It isn't important to them --- they just want it to work.
\item They are on a mission, and know they don't need to read everything.
\item If they find something that works, they stick to it rather than looking
      for a better way.
\end{enumerate}

\section{Billboard Design --- Design for Scanning, Not Reading}

\begin{itemize}
\item \sfb{Follow visual hierarchy} --- make important elements visually
      prominent, in proportion to their importance.
\item \sfb{Omit needless words} --- the art of not writing for the web. Happy
      talk must die; unnecessary instructions must die.
\item \sfb{Make clickables obvious} --- it must be instantly clear what can be
      clicked. Competition is one click away, and ambiguity produces frustration.
\item \sfb{Provide breadcrumbs and street signs} --- clear, conventional,
      intuitive navigation, so the user always knows where they are.
\end{itemize}

\section{Landing Page Essentials}

\begin{itemize}
\item \sfb{One page, one job.} A landing page with two goals achieves neither.
\item \sfb{A headline that matches the advertisement} that brought the visitor
      there. \emph{Message match} is the largest single lever on conversion rate.
\item \sfb{A clear, prominent, single call to action}, repeated above and below
      the fold.
\item \sfb{Benefit-led copy, not feature-led} --- what the visitor gets, not
      what the product has.
\item \sfb{Trust signals} --- testimonials, ratings, certifications, guarantees,
      secure-payment badges, real photographs.
\item \sfb{Minimal form friction} --- every additional field reduces completion.
      Ask only for what will actually be used.
\item \sfb{Fast load and mobile responsiveness} --- non-negotiable.
\item \sfb{Remove navigation and exits} on a dedicated campaign landing page, so
      the only way forward is the call to action\ix{call to action}.
\end{itemize}

\begin{summarybox}
Krug's first law is don't make me think: a page must be self-evident. Usability
is tested against seven attributes --- useful, learnable, memorable, effective,
efficient, desirable, delightful. People scan rather than read, satisfice rather
than optimise, and muddle through rather than learn, so pages must be designed
as billboards: visual hierarchy\ix{visual hierarchy}, no needless words, obvious clickables and
conventional navigation. A converting landing page has one job, a headline that
matches its advertisement, a single repeated call to action, benefit-led copy,
trust signals, minimal form friction, fast mobile load, and no exits.
\end{summarybox}

%%%%%%%%%%%%%%%%%%%%%%%%%%%%%%%%%%%%%%%%%%%%%%%%%%%%%%%%%%%%%%%%%%%%%%%%%%%%%%
\chapter{Measuring Impact --- Metrics and Performance}
%%%%%%%%%%%%%%%%%%%%%%%%%%%%%%%%%%%%%%%%%%%%%%%%%%%%%%%%%%%%%%%%%%%%%%%%%%%%%%

\syllabusstrip{Measure Impact; Metrics \& Performance.}

\begin{learningoutcomes}
\begin{itemize}
\item The three diagnostic buckets that give a metrics answer its structure.
\item The five metrics to name, and how each changes the next content cycle.
\item The rule that the metric must match the medium.
\item The five-step method for aligning KPIs to organisational goals.
\end{itemize}
\end{learningoutcomes}

\begin{keybox}[Structure Earns the Marks]
A list of metrics without a structure scores poorly. Organise any metrics answer
into the three buckets below, and for each metric state explicitly \emph{how the
reading changes the next decision}.
\pyq{CIE-I, 1 Jul 2024 --- 10 M; CIE-I, 16 Apr 2026 --- 10 M}
\end{keybox}

\section{The Three Diagnostic Buckets}

\begin{figure}[htbp]
\centering
\begin{tikzpicture}[font=\sffamily\footnotesize,
    bk/.style={rounded corners=4pt,line width=1pt,align=center,
               text width=3.35cm,inner sep=8pt,font=\sffamily\scriptsize}]
  \node[bk,draw=ocre,fill=ocrepale,text=slate] (b1) at (0,0)
    {\textbf{\textcolor{ocredark}{REACH \& VISIBILITY}}\\[4pt]
     impressions, unique viewers, reach, new users, backlinks earned\\[5pt]
     \textit{are we being seen at all?}};
  \node[bk,draw=deepblue,fill=deepbluepale,text=slate] (b2) at (4.05,0)
    {\textbf{\textcolor{deepblue}{RESONANCE \& MEMORY}}\\[4pt]
     brand recall, brand mentions, sentiment, share of voice, direct traffic
     growth\\[5pt]
     \textit{are we remembered?}};
  \node[bk,draw=greenok,fill=greenpale,text=slate] (b3) at (8.1,0)
    {\textbf{\textcolor{greenok}{ENGAGEMENT \& FRICTION}}\\[4pt]
     bounce rate, time on site, pages per session, scroll depth, shares,
     comments\\[5pt]
     \textit{is it actually good?}};
  \node[anchor=north,text=slatelight,font=\sffamily\scriptsize\itshape,
        align=center,text width=11.5cm] at (4.05,-2.6)
    {top-of-funnel scale \hfill lasting psychological impact \hfill content
     quality and user interaction};
\end{tikzpicture}
\caption[The three diagnostic buckets for content metrics]%
        {Each bucket diagnoses a different failure. High reach with low
         resonance is a memorability problem; high resonance with high friction
         is an execution problem.}
\label{fig:buckets}
\end{figure}

\section{The Five Metrics to Name}

\begin{enumerate}
\item \sfb{Website traffic and unique visitors} --- measures whether content is
      attracting an audience. \emph{Influence on strategy:} rising traffic on a
      theme signals which topics to double down on.
\item \sfb{Engagement rate} --- time on page, scroll depth\ix{scroll depth}, pages per session,
      likes, shares, comments; measures whether the content held attention.
      \emph{Influence:} a high-traffic, low-engagement page signals a headline
      that over-promises.
\item \sfb{Bounce rate} --- the percentage of visitors who leave after viewing
      only one page without interaction. \emph{Influence:} a high bounce\ix{bounce} on a
      converting keyword points to a message-match or a speed problem.
\item \sfb{Lead generation} --- form submissions, sign-ups, downloads; measures
      whether content converts attention into a contactable relationship.
      \emph{Influence:} identifies which assets deserve gated versions and paid
      promotion.
\item \sfb{Conversion rate and ROI} --- the percentage completing the desired
      action, and revenue attributable to content against its cost.
      \emph{Influence:} determines budget allocation for the next cycle.
\end{enumerate}

Further metrics worth naming if the question asks for more: organic search\ix{organic search}
ranking, keyword coverage, backlinks earned, e-mail open and click-through
rates, customer retention\ix{customer retention} and repeat purchase rate, and cost per lead.

\section{Align the Metric to the Medium}

\begin{warnbox}[The Measurement Rule]
Align the metric to the medium. Measuring an e-book by pageviews is a recipe for
failure.

\smallskip
Video \ra{} audience retention and total watch time. Blog and SEO \ra{} organic
rank and dwell time\ix{dwell time}. Lead magnets \ra{} download volume and lead conversion
rate. E-mail \ra{} open rate and click-through rate\ix{click-through rate}.
\end{warnbox}

\section{Aligning KPIs to Organisational Goals}

\begin{exambox}[Model Answer --- 10 Marks]
\pyq{CIE-I, 16 Apr 2026 --- 10 M}
\begin{enumerate}
\item \sfb{Identify business goals} --- define clear objectives: increase sales,
      brand awareness, lead generation\ix{lead generation}.
\item \sfb{Map goals to relevant KPIs} --- sales growth \ra{} conversion rate,
      revenue, cart-abandonment rate; brand awareness \ra{} traffic,
      impressions, reach; lead generation \ra{} form submissions, sign-ups;
      customer engagement \ra{} bounce rate, session duration, pages per visit.
\item \sfb{Make KPIs SMART} --- Specific, Measurable, Achievable, Relevant,
      Time-bound. For example: increase conversion rate by 10\% in three months.
\item \sfb{Use analytics tools} --- track through Google Analytics and
      dashboards.
\item \sfb{Monitor and optimise continuously} --- analyse regularly and modify
      strategy on the basis of measured performance.
\end{enumerate}
\end{exambox}

\begin{summarybox}
Structure every metrics answer into reach and visibility, resonance and memory,
and engagement and friction. Name five metrics --- traffic, engagement rate,
bounce rate, lead generation, and conversion rate with ROI --- and for each state
how the reading changes the next cycle. Match the metric to the medium: watch
time for video, rank and dwell for blogs, downloads for lead magnets, opens and
clicks for e-mail. Align KPIs to goals in five steps: identify goals, map to
KPIs, make them SMART\ix{SMART}, instrument with analytics, and optimise continuously.
\end{summarybox}

%%%%%%%%%%%%%%%%%%%%%%%%%%%%%%%%%%%%%%%%%%%%%%%%%%%%%%%%%%%%%%%%%%%%%%%%%%%%%%
\chapter{Using Content Research for Opportunities}
%%%%%%%%%%%%%%%%%%%%%%%%%%%%%%%%%%%%%%%%%%%%%%%%%%%%%%%%%%%%%%%%%%%%%%%%%%%%%%

\syllabusstrip{Using Content Research for Opportunities.}

\begin{learningoutcomes}
\begin{itemize}
\item The keyword sweet spot --- volume, difficulty and commercial intent\ix{commercial intent}.
\item The seven ways content research creates competitive edge.
\item The content research toolset, by job.
\end{itemize}
\end{learningoutcomes}

\begin{keybox}[High Yield]
This chapter answers a compulsory Question~2 sub-division in the Semester End
Examination.
\pyq{SEE, Jun/Jul 2025, Q2b --- 8 M}
\end{keybox}

\section{Finding the Keyword Sweet Spot}

A topic is worth pursuing only where three conditions overlap.

\begin{figure}[htbp]
\centering
\begin{tikzpicture}[font=\sffamily\footnotesize]
  \begin{scope}[opacity=0.40,transparency group]
    \fill[ocre]     (-1.3,0.6) circle (2.25);
    \fill[deepblue] ( 1.3,0.6) circle (2.25);
    \fill[greenok]  ( 0.0,-1.55) circle (2.25);
  \end{scope}
  \node[text=ocredark,font=\sffamily\bfseries\scriptsize,align=center]
    at (-2.55,1.75) {HIGH SEARCH\\ VOLUME};
  \node[text=deepblue,font=\sffamily\bfseries\scriptsize,align=center]
    at ( 2.55,1.75) {LOW SEO\\ DIFFICULTY};
  \node[text=greenok,font=\sffamily\bfseries\scriptsize,align=center]
    at ( 0.0,-3.4)  {COMMERCIAL INTENT};

  \node[text=slate,font=\sffamily\scriptsize,align=center,text width=2.1cm]
    at (-2.25,0.5) {roughly 250\\ to 1,000+\\ monthly};
  \node[text=slate,font=\sffamily\scriptsize,align=center,text width=2.1cm]
    at ( 2.25,0.5) {under about\\ 40 KD/SD};
  \node[text=slate,font=\sffamily\scriptsize,align=center,text width=2.4cm]
    at ( 0.0,-2.1) {high CPC, buyer-\\ specific phrasing};

  \node[text=white,font=\sffamily\bfseries\scriptsize,align=center]
    at (0,-0.32) {THE SWEET\\ SPOT};
\end{tikzpicture}
\caption[The keyword sweet spot]%
        {A topic is worth pursuing only in the central overlap. Two out of three
         is a trap: volume without intent is vanity, intent without
         achievability is wasted effort.}
\label{fig:sweetspot}
\end{figure}

\begin{itemize}
\item \sfb{High search volume} --- a minimum of roughly 250 to 1,000 or more
      monthly searches.
\item \sfb{Low SEO difficulty} --- under about a 40 difficulty score, so ranking
      is realistically achievable.
\item \sfb{Commercial intent} --- a high cost-per-click and buyer-specific
      phrasing, indicating the searcher intends to purchase.
\end{itemize}

\begin{warnbox}[Traffic Without Intent Is Useless]
Do not optimise for clicks; optimise for revenue. A page ranking first for a
query no buyer ever types is a vanity asset.
\end{warnbox}

\section{How Content Research Creates Competitive Edge}

\begin{exambox}[Model Answer --- 8 Marks]
\pyq{SEE, Jun/Jul 2025, Q2b --- 8 M}
\begin{enumerate}
\item \sfb{Reveals demand that already exists.} Keyword and question research
      shows what the market is actively asking, removing guesswork.
\item \sfb{Exposes content gaps} --- topics with volume where competitors rank
      poorly or not at all. This is the fastest route to visibility.
\item \sfb{Identifies trends early.} Google Trends and social listening surface
      rising topics before they saturate, allowing first-mover advantage.
\item \sfb{Benchmarks the competition.} It shows what already works in the
      niche, so effort goes to proven formats rather than experiments.
\item \sfb{Improves ROI.} Resources go to topics with demonstrated demand and
      achievable difficulty, instead of being spread evenly.
\item \sfb{Supplies the customer's own language.} Mining forums and reviews
      yields the exact phrasing to use in headlines and copy.
\item \sfb{Feeds every other channel.} One research pass supplies the SEO plan,
      the social calendar, the e-mail themes and the advertising copy.
\end{enumerate}
\end{exambox}

\section{The Content Research Toolset}

\begin{mmlongtable}{Content research tools by job}{tab:researchtools}%
{@{}P{0.22\textwidth}P{0.71\textwidth}@{}}
\rowcolor{ocre}\thd{Job} & \thd{Tools}\\
\midrule
\endfirsthead
\rowcolor{ocre}\thd{Job} & \thd{Tools}\\
\midrule
\endhead
\rlab{Keyword research} & Google Keyword Planner, SEMrush Keyword Magic Tool,
Ahrefs Keywords Explorer, Ubersuggest\\
\rowcolor{tablealt}
\rlab{Question mining} & AnswerThePublic, AlsoAsked, Google's ``People Also
Ask'' and ``Related searches''\\
\rlab{Trend detection} & Google Trends, Exploding Topics, platform trending
topics\\
\rowcolor{tablealt}
\rlab{Content performance} & BuzzSumo --- most-shared and most-linked content in
any topic; Ahrefs Content Explorer\\
\rlab{Audience listening} & Quora, Reddit, niche forums, marketplace and
app-store reviews, video comments\\
\rowcolor{tablealt}
\rlab{Gap analysis} & SEMrush Keyword Gap, Ahrefs Content Gap\\
\end{mmlongtable}

\begin{summarybox}
Pursue a topic only where high search volume, low SEO difficulty and commercial
intent overlap; two out of three is a trap. Content research creates edge by
revealing existing demand, exposing content gaps, catching trends early,
benchmarking competitors, improving ROI, supplying the customer's own language,
and feeding every other channel from a single pass. The toolset divides by job:
keyword research, question mining, trend detection, content performance,
audience listening and gap analysis.
\end{summarybox}

%%%%%%%%%%%%%%%%%%%%%%%%%%%%%%%%%%%%%%%%%%%%%%%%%%%%%%%%%%%%%%%%%%%%%%%%%%%%%%
\startappendices
%%%%%%%%%%%%%%%%%%%%%%%%%%%%%%%%%%%%%%%%%%%%%%%%%%%%%%%%%%%%%%%%%%%%%%%%%%%%%%

\chapter{Syllabus-to-Chapter Concordance}

This appendix exists so that coverage can be \emph{audited} rather than assumed.
The left-hand column reproduces the Unit~I syllabus descriptor phrase by phrase,
in the order the official document prints it. The right-hand column gives the
chapter and section that carries it. Every phrase resolves to a location, and no
numbered chapter of this volume contains material that is not traceable to a
phrase in this table.

\section*{First half --- Introduction to Digital Marketing}
\addcontentsline{toc}{section}{First half --- Introduction to Digital Marketing}

\begin{mmlongtable}{Concordance: Introduction to Digital Marketing}{tab:concordanceA}%
{@{}P{0.40\textwidth}P{0.53\textwidth}@{}}
\rowcolor{ocre}\thd{Syllabus phrase} & \thd{Where it is covered}\\
\midrule
\endfirsthead
\rowcolor{ocre}\thd{Syllabus phrase} & \thd{Where it is covered}\\
\midrule
\endhead
\rlab{Introduction to Digital Marketing} & Chapter~1 entire. Definition
\S\,1.4; traditional against digital \S\,1.5; push and pull \S\,1.6; the
marketing mix in the digital era \S\,1.7; enabling technology \S\,1.8; the
digital consumer \S\,1.9. Marketing foundations \S\,1.1--1.3.\\
\rowcolor{tablealt}
\rlab{Principles of Digital Marketing} & Chapter~2 entire. Five pillars
\S\,2.1; the internet paradox \S\,2.2; the seven governing principles \S\,2.3;
campaign characteristics \S\,2.4.\\
\rlab{Digital Marketing Channels} & Chapter~3 entire. Three senses of channel
\S\,3.1; paid, owned and earned \S\,3.2; hub and spokes \S\,3.3; affiliate
marketing \S\,3.4; channel choice \S\,3.5.\\
\rowcolor{tablealt}
\rlab{Tools to Create Buyer Persona} & Chapter~4 entire. Definition \S\,4.1;
five-step process \S\,4.2; factors \S\,4.3; canvas \S\,4.4; worked persona
\S\,4.5; toolset \S\,4.6.\\
\rlab{Competitor Research Tools} & Chapter~5 entire. What to look for \S\,5.1;
toolset \S\,5.2; converting findings into decisions \S\,5.3.\\
\rowcolor{tablealt}
\rlab{Website Analysis Tools} & Chapter~6 entire. Technical health against
behavioural insight \S\,6.1; toolset \S\,6.2; the six mechanisms \S\,6.3.\\
\end{mmlongtable}

\section*{Second half --- Content Marketing}
\addcontentsline{toc}{section}{Second half --- Content Marketing}

\begin{mmlongtable}{Concordance: Content Marketing}{tab:concordanceB}%
{@{}P{0.40\textwidth}P{0.53\textwidth}@{}}
\rowcolor{ocre}\thd{Syllabus phrase} & \thd{Where it is covered}\\
\midrule
\endfirsthead
\rowcolor{ocre}\thd{Syllabus phrase} & \thd{Where it is covered}\\
\midrule
\endhead
\rlab{Content Marketing Concepts \& Strategies} & Chapter~7 entire. Definition
\S\,7.1; a century of proof \S\,7.2; the compounding curve \S\,7.3; elite
traits \S\,7.4; the quality gate \S\,7.5.\\
\rowcolor{tablealt}
\rlab{Planning, Creating, Distributing \& Promoting Content} & Chapter~8
entire. The eight steps \S\,8.1; goal setting \S\,8.2; audience mapping
\S\,8.3; ideation and planning \S\,8.4; creation \S\,8.5; distribution
\S\,8.6; amplification \S\,8.7; the 5A path \S\,8.8; evaluation and improvement
\S\,8.9.\\
\rlab{Optimize Website UX \& Landing Pages} & Chapter~9 entire. Krug's first law
\S\,9.1; usability attributes \S\,9.2; how people use the web \S\,9.3;
billboard design \S\,9.4; landing page essentials \S\,9.5.\\
\rowcolor{tablealt}
\rlab{Measure Impact} & Chapter~10, \S\,10.1 and \S\,10.3 --- the diagnostic
buckets and the rule that the metric must match the medium.\\
\rlab{Metrics \& Performance} & Chapter~10, \S\,10.2 and \S\,10.4 --- the five
metrics and the alignment of KPIs to organisational goals.\\
\rowcolor{tablealt}
\rlab{Using Content Research for Opportunities} & Chapter~11 entire. The keyword
sweet spot \S\,11.1; the seven sources of edge \S\,11.2; the toolset \S\,11.3.\\
\end{mmlongtable}

\begin{keybox}[Reading the ``etc.'' in the Syllabus]
The official descriptor ends several clauses with ``etc.''. That word is treated
in this series as licence to complete a named framework, not as licence to add
new topics. So where the syllabus says ``Competitor Research Tools, Website
Analysis Tools, etc.'', the chapters cover those two tool families completely
rather than introducing a third family the syllabus never names.
\end{keybox}

\chapter{Previous-Year Question Bank --- Unit I}

Every question in this appendix is reproduced from an actual RVCE Marketing
Management Continuous Internal Evaluation or Semester End Examination paper for
this course. Answers and pointers follow the official schemes of valuation where
those were released. The bracketed tag after each question records the paper,
the date, and the mark allocation.

\section*{Part A --- Two-Mark Questions}
\addcontentsline{toc}{section}{Part A --- Two-Mark Questions}

\begin{enumerate}[leftmargin=2.2em,itemsep=1.1ex]

\item \textbf{Differentiate between marketplace and market space.}
      \pyqr{CIE-I, 3 Apr 2025; CIE-I, 16 Apr 2026 --- 2 M}\\
      Marketplace --- a physical or digital location where buyers and sellers
      come together to exchange goods and services (a mall, a showroom). Market
      space --- a conceptual or virtual environment where commercial value is
      created through digital networks and communication technologies (Amazon,
      an app). Shorthand: marketplace is physical, market space is digital.
      \hfill\emph{See} \S\,1.2.2.

\item \textbf{What is the mantra of marketing?}
      \pyqr{CIE-I, 3 Apr 2025 --- 2 M}\\
      Create, Communicate, Deliver, Value, Target, Profit.
      \hfill\emph{See} \S\,1.1.2.

\item \textbf{Define affiliate marketing and give an example.}
      \pyqr{CIE-I, 3 Apr 2025 --- 2 M}\\
      A performance-based marketing model in which a business rewards
      affiliates --- partners or individuals --- for driving traffic, leads or
      sales to its products or services through the affiliate's own marketing
      efforts. Example: Amazon Associates, where a blogger earns commission on
      purchases made via a tracked link. \hfill\emph{See} \S\,3.4.

\item \textbf{Mention any four characteristics of a good content.}
      \pyqr{CIE-I, 3 Apr 2025; CIE-I, 16 Apr 2026; SEE, Jun/Jul 2025 --- 2 M}\\
      Any four of: credible, shareable, useful, interesting, relevant,
      different, on-brand. The 2026 scheme also accepts relevant, accurate,
      engaging, clear, consistent, valuable. \hfill\emph{See} \S\,7.5.

\item \textbf{Differentiate between push and pull marketing.}
      \pyqr{CIE-I, 3 Apr 2025; CIE-I, 16 Apr 2026; SEE, Jun/Jul 2025 --- 2 M}\\
      Push --- the brand promotes products by pushing them towards the customer
      (television advertising, dealer incentives, direct selling). Pull --- the
      brand encourages customers to actively seek out the product (SEO, content,
      word of mouth\ix{word of mouth}). \hfill\emph{See} \S\,1.6.

\item \textbf{What is the first step of content marketing?}
      \pyqr{CIE-I, 16 Apr 2026 --- 2 M}\\
      Define goals and identify the target audience. \hfill\emph{See} \S\,8.2.

\item \textbf{Mention any two core concepts of marketing.}
      \pyqr{CIE-I, 16 Apr 2026 --- 2 M}\\
      Any two of: needs, wants, demand, value, exchange, markets.
      \hfill\emph{See} \S\,1.1.6.

\item \textbf{List any two principles of digital marketing.}
      \pyqr{SEE, Oct/Nov 2023 --- 2 M}\\
      Any two of the five pillars --- know your business; know the competition;
      know your customers; know what you want to achieve; know how you're doing.
      Alternatively: measurability, interactivity, personalisation,
      permission-based targeting. \hfill\emph{See} \S\,2.1 and \S\,2.3.

\item \textbf{Differentiate between market segmentation and market mix.}
      \pyqr{SEE, Oct/Nov 2023 --- 2 M}\\
      Market segmentation is the process of dividing a heterogeneous total
      market into smaller homogeneous groups of buyers with similar needs, on
      bases such as geographic, demographic, psychographic and behavioural, so
      that each can be targeted distinctly. The marketing mix is the set of
      controllable tactical tools --- Product, Price, Place and Promotion ---
      that the firm blends to produce the response it wants from its chosen
      target segment. In short: segmentation decides \emph{whom} to serve; the
      marketing mix decides \emph{how} to serve them.
      \hfill\emph{See} \S\,1.7 and Appendix~C.3.

\end{enumerate}

\section*{Part B --- Eight- and Ten-Mark Questions}
\addcontentsline{toc}{section}{Part B --- Eight- and Ten-Mark Questions}

\begin{enumerate}[leftmargin=2.2em,itemsep=1.4ex]

\item \textbf{A newly launched e-commerce company wants to build brand awareness
      through content marketing. Suggest a step-by-step strategy for them,
      including content types and distribution channels.}
      \pyqr{CIE-I, 3 Apr 2025 --- 10 M}\\
      Official scheme answer --- the eight steps: goal setting, audience
      mapping, content ideation and planning, content creation, content
      distribution, content amplification, content marketing evaluation, content
      marketing improvement. Expand each with the detail in Chapter~8, naming
      content types (blogs, videos, infographics, podcasts, e-books,
      user-generated content) and distribution across owned, earned and paid
      media. \hfill\emph{See} \S\,8.1 and \S\,3.2.

\item \textbf{A newly launched health \& wellness startup offering organic
      lifestyle products aims to build strong brand awareness through content
      marketing. Design a step-by-step content marketing strategy, clearly
      outlining: content formats; content themes; distribution platforms (owned,
      earned, paid media\ix{paid media}); metrics for performance evaluation.}
      \pyqr{CIE-I, 16 Apr 2026 --- 10 M}\\
      The full five-part scheme answer is set out in Chapter~8 as a worked
      scenario. \hfill\emph{See} \S\,8.9.

\item \textbf{A travel agency specializing in adventure tourism wants to improve
      its targeted marketing. How should they create a buyer persona, and what
      factors (e.g., demographics, interests, online behaviour) should they
      consider?}
      \pyqr{CIE-I, 3 Apr 2025; CIE-I, 16 Apr 2026 --- 10 M}\\
      Answer in three parts. (i) \emph{Steps:} collect data through surveys,
      website analytics and social media insights; segment the audience by
      shared traits; identify goals and pain points; build the persona profile
      with a name, background and behaviour details; validate and update using
      feedback and trends. (ii) \emph{Factors:} demographics (age 20--40,
      gender, income, occupation, location); psychographics (trekking, scuba
      diving, skydiving, fitness and outdoor lifestyle, risk-taking attitude,
      preference for unique experiences over luxury); online behaviour
      (Instagram, YouTube, travel blogs; searches such as ``best trekking
      places''; prefers videos and reviews; books online via mobile apps);
      travel behaviour (frequency, budget against premium, group against solo,
      preferred destinations); pain points (safety concerns, budget constraints,
      lack of information, trust in agencies). (iii) \emph{Example persona:}
      Adventure Arun, 28, IT professional; interests trekking and camping;
      follows travel influencers and books online; goal is unique thrilling
      experiences; pain point is safety and affordability.
      \hfill\emph{See} Chapter~4.

\item \textbf{A retail company has been using traditional marketing methods such
      as newspaper ads, billboards, and TV commercials. However, they are now
      considering a shift to digital marketing. Compare the benefits and
      limitations of both approaches and suggest a strategy for transitioning to
      digital marketing.}
      \pyqr{CIE-I, 3 Apr 2025; CIE-I, 16 Apr 2026 --- 10 M}\\
      \emph{Traditional benefits:} wide reach through television and newspapers;
      high credibility and trust; effective for a mass audience.
      \emph{Traditional limitations:} expensive; difficult to measure results;
      no precise targeting; one-way communication. \emph{Digital benefits:}
      cost-effective; targeted by age, location and interest; measurable through
      analytics and ROI; two-way interaction; global reach. \emph{Digital
      limitations:} requires technical skills; high online competition; privacy
      concerns; dependent on internet access. \emph{Transition strategy:} the
      eight sequenced steps in \S\,1.5.3.
      \hfill\emph{See} \S\,1.5.

\item \textbf{Identify and explain the various entities that can be marketed,
      providing relevant examples for each category.}
      \pyqr{CIE-I, 3 Apr 2025 --- 10 M}\\
      The ten entities --- goods, services, experiences, events, persons, places,
      properties, organizations, information and ideas --- each explained with an
      example. \hfill\emph{See} \S\,1.1.5, Table~\ref{tab:entities}.

\item \textbf{How have the 4Ps of marketing (Product, Price, Place, and
      Promotion) evolved in the digital era? Provide examples of companies
      effectively leveraging digital strategies for each element.}
      \pyqr{CIE-I, 3 Apr 2025 --- 10 M}\\
      Product \ra{} digital and personalised products (Spotify's data-driven
      playlists). Price \ra{} dynamic and transparent pricing (Amazon's
      AI-driven adjustment by demand, competition and user behaviour). Place
      \ra{} online channels and omnichannel presence (Nike's seamless website,
      app and store experience with click-and-collect). Promotion \ra{} digital
      and interactive marketing (Coca-Cola's ``Share a Coke'' campaign).
      \hfill\emph{See} \S\,1.7.2, Table~\ref{tab:digitalps}.

\item \textbf{Describe the key principles of digital marketing and explain how
      they differ from traditional marketing approaches. Provide examples to
      illustrate your points.}
      \pyqr{CIE-I, 1 Jul 2024 --- 10 M}\\
      State the five pillars, add measurability, interactivity,
      personalisation, permission and iteration, then contrast each against
      traditional marketing using Table~\ref{tab:tradvsdig}, with one concrete
      example per point. \hfill\emph{See} Chapter~2 and \S\,1.5.2.

\item \textbf{Your company has a limited marketing budget and needs to choose
      between Search Engine Optimization, Pay Per Click, and Social Media
      Marketing to promote a new product. Based on the product being a high-tech
      gadget aimed at young professionals, which channel(s) would you prioritize
      and why? Justify your decision with specific examples.}
      \pyqr{CIE-I, 1 Jul 2024 --- 10 M}\\
      Compare the three on speed, cost model, longevity, buyer intent and
      predictability, then commit to a weighted recommendation and justify it
      with the intent argument --- a brand-new gadget generates no organic search
      demand yet, so paid and social must create the demand that SEO will later
      harvest. \hfill\emph{See} \S\,3.5.

\item \textbf{Identify and explain at least five key metrics used to measure the
      impact and performance of content marketing efforts. How can these metrics
      influence future content strategies?}
      \pyqr{CIE-I, 1 Jul 2024 --- 10 M}\\
      Name traffic, engagement rate, bounce rate, lead generation and conversion
      rate with ROI; group them into reach and visibility, resonance and memory,
      and engagement and friction; and for each state explicitly how the reading
      changes the next content cycle. \hfill\emph{See} Chapter~10.

\item \textbf{Imagine you are a digital marketing strategist at GreenTech
      Innovations, a company launching a new line of eco-friendly kitchen
      appliances called ``EcoChef.'' Design a comprehensive content marketing
      strategy to promote these products. Include specific goals, types of
      content, distribution channels, and methods of engagement with your target
      audience.}
      \pyqr{CIE-I, 1 Jul 2024 --- 10 M}\\
      The full worked answer is set out in Chapter~8.
      \hfill\emph{See} \S\,8.9.

\item \textbf{What is Digital marketing? Explain the differences between
      Traditional Marketing and Digital Marketing. Highlight any 5 key
      differences.}
      \pyqr{SEE, Oct/Nov 2023, Q2a --- 8 M}\\
      Define digital marketing, then give any five differences from
      Table~\ref{tab:tradvsdig} --- communication direction, reach, cost,
      measurability and interactivity are the strongest five --- each with a
      one-line example. \hfill\emph{See} \S\,1.4 and \S\,1.5.2.

\item \textbf{List and explain various pricing strategies used in marketing
      management with suitable examples.}
      \pyqr{SEE, Oct/Nov 2023, Q2b --- 8 M}\\
      Ten strategies with examples are set out in Appendix~C.2, which also
      records why this topic sits outside the literal Unit~I descriptor.
      \hfill\emph{See} Appendix~C.2.

\item \textbf{Explain how content research tools help marketers discover new
      opportunities and gain a competitive edge.}
      \pyqr{SEE, Jun/Jul 2025, Q2b --- 8 M}\\
      The keyword sweet spot (volume $\times$ low difficulty $\times$ commercial
      intent), the seven ways research creates edge, and the toolset.
      \hfill\emph{See} Chapter~11.

\item \textbf{Discuss how website analysis contributes to enhanced customer
      experience and conversion optimization.}
      \pyqr{SEE, Jun/Jul 2025, Q2a --- 8 M}\\
      Name the tools by category, then give the six mechanisms: diagnosing
      friction, fixing speed, locating funnel drop-off, validating changes
      through A/B testing, segmenting by device and channel, and aligning
      content with intent. \hfill\emph{See} Chapter~6.

\item \textbf{Enumerate briefly on the best practices for creating compelling and
      shareable content on social media platforms to increase brand visibility
      and engagement.}
      \pyqr{SEE, Oct/Nov 2023, Q3b --- 8 M}\\
      Apply the seven-point quality gate; lead with visuals (visual content is
      far more likely to be shared, and posts with relevant images earn
      substantially more views); follow the 4-1-1 rule so that value outweighs
      promotion; write for the platform rather than cross-posting blindly; hook
      in the first three seconds; use one or two targeted hashtags; post
      consistently at audience-active times; invite participation through
      questions, polls and contests; repurpose one master asset into many
      formats; and amplify through communities and influencers.
      \hfill\emph{See} \S\,7.5 and \S\,8.7. \emph{The 4-1-1 rule and the
      platform-specific practices belong to Unit~II; this question is
      set from Unit~I in the paper but answered across both.}

\end{enumerate}

\section*{How Unit I Is Examined}
\addcontentsline{toc}{section}{How Unit I Is Examined}

\begin{keybox}[The Pattern Across Papers]
In the Semester End Examination, Question~2 is compulsory and is always drawn
from Unit~I --- two sub-divisions of eight marks each --- plus roughly two Part~A
objective questions. In the first Continuous Internal Evaluation test, Unit~I
typically carries the whole paper.

\smallskip
The three highest-probability ten-mark topics are the eight-step content
marketing strategy applied to a scenario, buyer persona creation, and
traditional against digital marketing with a transition plan. All three have
already repeated across multiple papers.
\end{keybox}

\chapter{Supplementary Examined Topics}

\begin{warnbox}[What This Appendix Is, and Why It Is Not a Chapter]
The three topics below appear in past papers for this course but sit outside the
literal wording of the Unit~I syllabus descriptor --- and, in fact, outside the
wording of any of the five descriptors. They are placed in an appendix rather
than in a numbered chapter so that the rule governing this series is not broken:
\emph{the numbered chapters carry the syllabus and nothing else}. They are
included rather than dropped because they are genuinely examined, and a student
who skipped them would lose real marks.
\end{warnbox}

\section{The Product Life Cycle}

Every product follows a largely inevitable path from launch to eventual
disappearance, and the stage a product occupies constrains its content and
channel decisions --- which is why the topic appears in the supplementary
lecture material for this unit.

\subsection{Diagnosing the stage}

\begin{itemize}
\item \sfb{Market volatility} --- are rapid industry shifts limiting the
      product's lifespan, or favouring new and improved alternatives?
\item \sfb{Consumer behaviour} --- is the product an everyday necessity? Does it
      solve a common problem? Is its usage simple or complex?
\item \sfb{Competitive landscape} --- how easily can competitors replicate it?
      Map the number of rivals, their expertise, pricing and distribution
      channels.
\end{itemize}

\begin{figure}[htbp]
\centering
\begin{tikzpicture}
\begin{axis}[
  width=12.2cm, height=5.9cm,
  font=\sffamily\footnotesize,
  xlabel={time}, ylabel={sales volume},
  xlabel style={font=\sffamily\scriptsize,text=slate},
  ylabel style={font=\sffamily\scriptsize,text=slate},
  axis lines=left, axis line style={rulegrey},
  xtick=\empty, ytick=\empty,
  xmin=0, xmax=10.4, ymin=0, ymax=1.34,
  clip=false]

  % Plotted from explicit coordinates rather than a closed-form expression:
  % the classic curve is piecewise (slow start, steep growth, long plateau,
  % decline) and no single smooth formula gives all four stages the right
  % relative width.
  \addplot[ocre,very thick,smooth,tension=0.55] coordinates {
    (0.00,0.02) (0.70,0.05) (1.35,0.11)
    (2.10,0.26) (2.90,0.48) (3.90,0.72)
    (4.80,0.87) (5.60,0.95) (6.40,0.98)
    (7.10,0.97) (7.60,0.94) (8.30,0.83)
    (9.00,0.66) (9.70,0.46) (10.20,0.33)};

  % Stage boundaries, written out rather than looped: pgfplots defers parsing
  % of "axis cs:" coordinates until \end{axis}, by which time a \foreach loop
  % variable no longer exists.
  \draw[rulegrey,dashed] (axis cs:1.35,0) -- (axis cs:1.35,1.20);
  \draw[rulegrey,dashed] (axis cs:3.90,0) -- (axis cs:3.90,1.20);
  \draw[rulegrey,dashed] (axis cs:7.60,0) -- (axis cs:7.60,1.20);

  % stage names, centred in each band
  \node[text=ocredark,font=\sffamily\bfseries\scriptsize] at (axis cs:0.70,1.26) {INTRODUCTION};
  \node[text=ocredark,font=\sffamily\bfseries\scriptsize] at (axis cs:2.63,1.26) {GROWTH};
  \node[text=ocredark,font=\sffamily\bfseries\scriptsize] at (axis cs:5.75,1.26) {MATURITY};
  \node[text=ocredark,font=\sffamily\bfseries\scriptsize] at (axis cs:9.10,1.26) {DECLINE};
\end{axis}
\end{tikzpicture}
\caption[The product life cycle]%
        {The product life cycle\ix{product life cycle}. Maturity is normally the longest stage, and the
         curve's shape is not fixed --- see \S\,C.1.3.}
\label{fig:plc}
\end{figure}

\subsection{The master lifecycle matrix}

\begin{mmlongtable}{The product life cycle matrix}{tab:plcmatrix}%
{@{}P{0.147\textwidth}P{0.187\textwidth}P{0.187\textwidth}P{0.187\textwidth}P{0.177\textwidth}@{}}
\rowcolor{ocre}\thd{} & \thd{Introduction} & \thd{Growth} & \thd{Maturity} & \thd{Decline}\\
\midrule
\endfirsthead
\rowcolor{ocre}\thd{} & \thd{Introduction} & \thd{Growth} & \thd{Maturity} & \thd{Decline}\\
\midrule
\endhead
\rlab{Sales volume} & Low & Surging & Stabilising & Dropping\\
\rowcolor{tablealt}
\rlab{Competition} & Non-existent & Increasing & Entrenched & Dominating\\
\rlab{Primary goal} & Create awareness and generate demand & Build market share
and loyalty & Defend position and maximise profit & Minimise losses\\
\rowcolor{tablealt}
\rlab{Core strategy} & Promotions and demonstrations & Customisation and
expansion & Diversification and cross-selling & Niche focus and liquidation\\
\rlab{Historical example} & Apple iPhone (2007) & Zara (1990s--2000s) &
Kellogg's & Kodak (1990s--2012)\\
\end{mmlongtable}

\subsection{Stage tactics, and why the curve is not always linear}

\begin{description}
\item[Introduction] Consumers are unfamiliar; demand is low while production and
distribution costs are high. Offer targeted promotions, discounts and free
samples; host product demonstrations to prove functional value; deploy influencer
marketing for immediate social proof; use mass media advertising; create
shareable brand experiences. \emph{Apple, 2007:} massive television advertising
plus a limited initial release and heavily promoted exclusivity generated
unprecedented demand.
\item[Growth] Sales and profits surge as competitors flood the market. Introduce
product customisation; deliver exceptional customer service to lock in loyalty;
identify new segments and expand geographically. \emph{Zara:} rapid
international expansion while producing small quantities of exclusive designs,
creating permanent novelty.
\item[Maturity] The longest stage; growth slows, competitors are entrenched and
prices fall. Diversify product lines; drive cross-selling with complementary
goods; innovate continuously through formula or feature updates; deploy
value-based content marketing; compete on cost by optimising production and
sourcing. \emph{Kellogg's:} competitor-comparison advertising, organic lines,
and cereal bundled with dried fruit and snack bars.
\item[Decline] Demand drops as technology or trends shift. Focus on a smaller,
highly profitable niche; reposition the brand for new segments; cut marketing
spend and rely on organic channels; implement deep cost reductions and
discounting; execute exit strategies. \emph{Kodak:} cut prices, diversified into
digital printing and sold patents, yet still filed for bankruptcy --- proving
that some technological shifts are irreversible.
\end{description}

Three deviations from the standard curve are worth naming.

\begin{itemize}
\item \sfb{Short life} --- fad products that spike and die without ever reaching
      maturity.
\item \sfb{Extended life} --- core staples that plateau and dominate profitably
      for decades.
\item \sfb{The pivot} --- products that reset their own curve through radical
      repositioning or new-market expansion.
\end{itemize}

\section{Pricing Strategies}

\begin{exambox}[Model Answer --- 8 Marks]
\pyq{SEE, Oct/Nov 2023, Q2b --- 8 M}
Ten strategies, each with an example. Naming and explaining any six or seven
with examples is normally sufficient at eight marks.
\end{exambox}

\begin{mmlongtable}{Pricing strategies}{tab:pricing}%
{@{}P{0.22\textwidth}P{0.42\textwidth}P{0.29\textwidth}@{}}
\rowcolor{ocre}\thd{Strategy} & \thd{Basis} & \thd{Example}\\
\midrule
\endfirsthead
\rowcolor{ocre}\thd{Strategy} & \thd{Basis} & \thd{Example}\\
\midrule
\endhead
\rlab{Penetration pricing} & A low launch price to win share fast & Jio's initial
free and low tariffs\\
\rowcolor{tablealt}
\rlab{Skimming pricing} & A high launch price harvested from early adopters, then
lowered & New smartphone flagship models\\
\rlab{Cost-plus pricing} & Cost plus a fixed markup & Retail and FMCG
distribution\\
\rowcolor{tablealt}
\rlab{Value-based pricing} & Priced on perceived customer value rather than cost
& Premium software, luxury goods\\
\rlab{Competitive or going-rate pricing} & Matched to rivals & Airlines, fuel\\
\rowcolor{tablealt}
\rlab{Psychological pricing} & Prices set just below a round number &
\rupee 999 instead of \rupee 1,000\\
\rlab{Bundle pricing} & Several items priced together below their sum &
Combo packs, office software suites\\
\rowcolor{tablealt}
\rlab{Freemium} & A free core product with a paid upgrade & Streaming and
conferencing services\\
\rlab{Dynamic pricing} & Real-time adjustment by demand & Ride-hailing surge
pricing, large marketplaces\\
\rowcolor{tablealt}
\rlab{Promotional or discount pricing} & Temporary reductions to drive volume &
Festival mega-sale events\\
\end{mmlongtable}

\section{Market Segmentation Against the Marketing Mix}

\begin{definitionbox}[Market Segmentation]
\keytermas{Market segmentation}{market segmentation} is the process of dividing a heterogeneous total
market into smaller homogeneous groups of buyers with similar needs, on bases
such as geographic, demographic, psychographic and behavioural, so that each
group can be targeted distinctly.
\end{definitionbox}

The distinction the examiner wants is one of \emph{function}, and it reduces to
one line.

\begin{keybox}[The One-Line Distinction]
\textbf{Segmentation decides whom to serve. The marketing mix decides how to
serve them.}
\pyq{SEE, Oct/Nov 2023 --- 2 M}

\smallskip
Segmentation is an analytical, prior decision about the market; the mix is a set
of controllable tactical levers applied after that decision has been taken.
Changing the segment changes every element of the mix, but changing the mix does
not change the segment.
\end{keybox}

\chapter{Glossary of Key Terms}

\begin{mmlongtable}{Key terms for rapid revision}{tab:glossary}%
{@{}P{0.26\textwidth}P{0.67\textwidth}@{}}
\rowcolor{ocre}\thd{Term} & \thd{One-line meaning}\\
\midrule
\endfirsthead
\rowcolor{ocre}\thd{Term} & \thd{One-line meaning}\\
\midrule
\endhead
\rlab{Affiliate marketing} & Performance-based marketing paying partners for
traffic, leads or sales actually delivered\\
\rowcolor{tablealt}
\rlab{Bounce rate} & The percentage of visitors who leave after viewing only one
page without interaction\\
\rlab{Brand} & An offering from a known source\\
\rowcolor{tablealt}
\rlab{Buyer persona} & A research-based, semi-fictional profile of the ideal
customer\\
\rlab{Call to action (CTA)} & The explicit instruction telling the visitor what
to do next\\
\rowcolor{tablealt}
\rlab{Content marketing} & Creating and distributing valuable, relevant,
consistent content to drive profitable customer action\\
\rlab{Demands} & Wants for a specific product backed by ability and willingness
to pay\\
\rowcolor{tablealt}
\rlab{Digital marketing} & Targeted, conversion-oriented, quantifiable,
interactive marketing using digital technologies\\
\rlab{Earned media} & Reach the audience gives you and you did not pay for\\
\rowcolor{tablealt}
\rlab{Holistic marketing} & The view that everything matters in marketing:
internal, integrated, performance and relationship marketing\\
\rlab{Influencer} & An early-adopter online opinion leader who holds the virtual
ear of the masses\\
\rowcolor{tablealt}
\rlab{Landing page} & A single-purpose page a campaign sends traffic to, built
for one conversion\\
\rlab{Mantra of marketing} & Create, Communicate, Deliver, Value, Target, Profit\\
\rowcolor{tablealt}
\rlab{Marketing management} & Choosing target markets and getting, keeping and
growing customers through superior customer value\\
\rlab{Marketing mix} & The controllable tactical tools: Product, Price, Place,
Promotion\\
\rowcolor{tablealt}
\rlab{Marketplace} & A physical or digital location where exchange occurs\\
\rlab{Market space} & A virtual environment in which value is created digitally\\
\rowcolor{tablealt}
\rlab{Message match} & The correspondence between an advertisement and the
landing page it leads to\\
\rlab{Metamarket} & A cluster of complementary products and services related in
the consumer's mind but spread across industries\\
\rowcolor{tablealt}
\rlab{Needs} & Basic human requirements, not created by marketers\\
\rlab{Offering} & The combination of products, services, information and
experiences that delivers the value proposition\\
\rowcolor{tablealt}
\rlab{Owned media} & Assets entirely controlled by the brand\\
\rlab{Paid media} & Reach bought with capital, which ends when spending ends\\
\rowcolor{tablealt}
\rlab{Perceived anonymity} & The online trait that frees consumers from
real-world social constraint\\
\rlab{Product life cycle} & Introduction, growth, maturity, decline\\
\rowcolor{tablealt}
\rlab{Prosumer} & A consumer who also shapes or produces what they consume\\
\rlab{Pull marketing} & The customer actively seeks the brand out\\
\rowcolor{tablealt}
\rlab{Push marketing} & The brand pushes the product towards the customer\\
\rlab{QSP triad} & Quality, service and price --- the components of value\\
\rowcolor{tablealt}
\rlab{Satisficing} & Choosing the first option that plausibly works rather than
the best one\\
\rlab{Satisfaction} & A person's judgement of perceived performance against
expectations\\
\rowcolor{tablealt}
\rlab{Share of voice} & Conversation about your brand as a proportion of
conversation about the category\\
\rlab{Value proposition} & The set of benefits a firm offers to satisfy a
customer need\\
\rowcolor{tablealt}
\rlab{Wants} & Needs directed at specific objects, shaped by society and
personality\\
\end{mmlongtable}

\chapter{Framework Quick-Reference Sheet}

Everything in this volume that is a numbered list, in one place, for the night
before the examination.

\begin{mmlongtable}{Framework quick reference}{tab:quickref}%
{@{}P{0.27\textwidth}P{0.66\textwidth}@{}}
\rowcolor{ocre}\thd{Framework} & \thd{The sequence}\\
\midrule
\endfirsthead
\rowcolor{ocre}\thd{Framework} & \thd{The sequence}\\
\midrule
\endhead
\rlab{Mantra of marketing} & Create \ra{} Communicate \ra{} Deliver \ra{} Value
\ra{} Target \ra{} Profit\\
\rowcolor{tablealt}
\rlab{Ten marketable entities} & Goods, services, events, experiences, persons,
places, properties, organizations, information, ideas\\
\rlab{Core concepts} & Needs, wants, demands, value, exchange, markets\\
\rowcolor{tablealt}
\rlab{Value creation engine} & Audience \ra{} solution \ra{} result\\
\rlab{Satisfaction rule} & Performance below expectation dissatisfies; equal
satisfies; above delights\\
\rowcolor{tablealt}
\rlab{Five company orientations} & Production, product, selling, marketing,
holistic\\
\rlab{Holistic marketing} & Internal, integrated, performance, relationship\\
\rowcolor{tablealt}
\rlab{Eight management tasks} & Strategies and plans, insights, connecting,
brands, creating, delivering, communicating, long-term growth\\
\rlab{Marketing mix} & Product, Price, Place, Promotion (+ People, Process,
Physical evidence; mirrored by the four Cs)\\
\rowcolor{tablealt}
\rlab{Technology adoption pattern} & Emergence \ra{} foothold \ra{} exploration
\ra{} mainstream\\
\rlab{Five online consumer traits} & Comfortable, want it now, in control,
fickle, vocal\\
\rowcolor{tablealt}
\rlab{Five strategy pillars} & Know your business, the competition, your
customers, your goals, your results\\
\rlab{Seven governing principles} & Measurability, interactivity,
personalisation, permission, integration, iteration, content centrality\\
\rowcolor{tablealt}
\rlab{Media taxonomy} & Owned (foundation), paid (accelerator), earned
(validator)\\
\rlab{Persona process} & Research customers \ra{} segment \ra{} gather data
\ra{} identify goals and pain points \ra{} build and validate\\
\rowcolor{tablealt}
\rlab{Persona canvas} & Background and goals; information sources; messaging and
topics; purchase role and objections\\
\rlab{Competitor decisions} & Emulate, avoid, occupy\\
\rowcolor{tablealt}
\rlab{Website analysis halves} & Technical health; behavioural insight\\
\rlab{Content quality gate} & Credible, shareable, useful or fun, interesting,
relevant, different, on-brand\\
\rowcolor{tablealt}
\rlab{Elite content marketer traits} & Documented strategy, dedicated
leadership, expanded tactics, multi-platform distribution, budget allocation\\
\rlab{Eight content steps} & Goals \ra{} audience \ra{} ideation \ra{} creation
\ra{} distribution \ra{} amplification \ra{} evaluation \ra{} improvement\\
\rowcolor{tablealt}
\rlab{Three content phases} & Architecture, mechanics, telemetry\\
\rlab{Audience mapping layers} & Anxieties and desires \ra{} pain points \ra{}
profile and persona\\
\rowcolor{tablealt}
\rlab{Content plan layers} & Strategic, structural, temporal \ra{} road map\\
\rlab{Amplification nodes} & Communities, buzzers, influencers\\
\rowcolor{tablealt}
\rlab{5A customer path} & Aware, appeal, ask, act, advocate --- visible,
relatable, searchable, actionable, shareable\\
\rlab{Content flywheel} & Audience research \ra{} keyword strategy \ra{}
creation \ra{} distribution \ra{} measurement \ra{} repeat\\
\rowcolor{tablealt}
\rlab{Seven usability attributes} & Useful, learnable, memorable, effective,
efficient, desirable, delightful\\
\rlab{Billboard design rules} & Visual hierarchy, omit needless words, obvious
clickables, breadcrumbs and street signs\\
\rowcolor{tablealt}
\rlab{Three metric buckets} & Reach and visibility; resonance and memory;
engagement and friction\\
\rlab{Five content metrics} & Traffic, engagement rate, bounce rate, lead
generation, conversion rate and ROI\\
\rowcolor{tablealt}
\rlab{KPI alignment} & Identify goals \ra{} map to KPIs \ra{} make SMART \ra{}
instrument analytics \ra{} optimise continuously\\
\rlab{Keyword sweet spot} & High volume $\times$ low difficulty $\times$
commercial intent\\
\rowcolor{tablealt}
\rlab{Product life cycle} & Introduction, growth, maturity, decline\\
\end{mmlongtable}

%%%%%%%%%%%%%%%%%%%%%%%%%%%%%%%%%%%%%%%%%%%%%%%%%%%%%%%%%%%%%%%%%%%%%%%%%%%%%%
\backmatter
\printindex
\end{document}
